% Modernised reading — fonds Alexandre Grothendieck, shelfmark 72
% (pages 1-112, the whole folder).
%
% This is an INTERPRETATION, not a transcription. It re-reads the pages in
% today's notation and vocabulary, states the hypotheses the manuscript leaves
% implicit, and drops the critical apparatus so that the mathematics reads
% straight through. Everything the brackets and underlines carried moves into
% a footnote. It is held to being mathematically correct as it stands; where
% the manuscript is loose or mistaken, the true statement is what appears here
% and the footnote says what the page has. In any dispute of substance the
% transcription governs: whoever cites Grothendieck must cite batch-01.fr.tex
% (pp. 1-20), batch-02 (21-40), batch-03 (41-60), batch-04 (61-80), batch-05
% (81-100) or batch-06.fr.tex (101-112).
%
% Pass: Opus 5.5 (claude-opus-5-5), 2026-10-10 — first pass, unchecked against the pages by a human.
% The folder was read whole, its six batches together; this is one document
% for the shelfmark. It was made on Opus 5.5 and has not been measured against
% the reference reading of folder 115 (made on Opus 5). The literature named
% in the footnotes is cited from memory and was not checked against the
% sources.
%
% Scope. 112 pages, 93 of them transcribed with a \page marker. No marker:
% the wrappers 1, 26, 31, 36, 43, 55, 71, 78, 83, 95, 106 (titles in his hand,
% which become section titles here) and 21 (blank); the typed USTL documents
% 59, 61, 63, 65, 67, 69 (avis de soutenance of June-July 1978, a « Recueil »
% of Conseil de l'Université decisions of May 1978), with nothing of his:
% named, not described. Page 30 carries no marker and is not accounted for in
% the batch-2 header (which says « 27-30 » for the hermitian sheets, the wrapper
% of page 26 says « (4p) »); this reading cannot say what it holds.
%
% Spine. The linear-algebraic theory of geometric realisations of regular
% polytopes by (pseudo-)reflections, over an arbitrary base scheme S:
%   2-20    the « dictionnaire »: a string of equivalent categories (projective,
%           affine, vector, at infinity, orthogonal), the antipodism (first
%           state), duality, twisting by torsors, the invariant quadratic form.
%   22-42   antipodism (second state), the group scheme U_rho, the hermitian
%           variant (27-29), the reconstruction of tau_i from rotations (32-35).
%   44-60   discriminants: the Gram matrix, continuants L_k, the centre c~,
%           rho, sigma, smoothness, the antipodal plane; dual bases (60).
%   62-77   the dual system and the centred computations: circumsphere,
%           midsphere, discriminant delta(f_V) = B rho.
%   79-82   invariant subspaces over a field.
%   84-105  foundations: realisations of a regular n-carte, non-degeneracy,
%           the torsor of flags (G_m or G_a), universal coordinates, the lifting
%           lemma, the Théorème (A) <=> (B) and the existence of s_0.
%   107-112 Stiefel-Whitney classes and Brauer class of the invariant form.
% The order of the leaves is the archivists'; the foundations (84-105) are
% logically first and are cited as such.
%
% Notation fixed for the whole document (his letters move between sheets):
%   pi in E^vee the affine form (his pi on pp. 22-25, his pi' on pp. 5-9 and
%   37-112); c~ in E the centre vector (his sigma~ / c~ / c' on pp. 22, 47-53,
%   his pi on pp. 5-9, 38-40, 60-70, 107; looped sigma on p. 77); c = c~/rho,
%   c' = 2c. rho~, sigma~ = the « halved » continuants L' (his rho', sigma' on
%   pp. 22-24, 52-54, 74-77; his tildes on pp. 39-42, 62-66, 107); sigma^vee =
%   L_n(beta_0..beta_{n-2}) and gamma = L_{n-1}(beta_1..beta_{n-2}) (his sigma',
%   gamma of pp. 60-70). nu_i the products of the Théorème of p. 96 (his lambda_i
%   there, not the eigenvalues). f_E the invariant form on E (his f, f^eps,
%   f_eps, the subscript read eps being script E).
%
% Substantive departures from the pages, each footnoted where it occurs:
%   p. 13 margin: 2n-1 parameters read 2n+1 (the list he gives has 2n+1);
%   p. 16: « si n >= 2 (mais non si n = 2) » — the equivalence holds for n >= 2
%     and fails for n = 1 (lifting lemma of p. 94 needs rank >= 3);
%   p. 17: cross terms of f_V carry a minus sign (pp. 44, 60 agree);
%   p. 18: c') is (tau_i tau_{i+1})^2 != id, the page writes = id;
%   p. 18: b) => b') is not true in general (theorem of pp. 79-82);
%   p. 35: names alpha/beta swapped and 2(2cos^2+1) = 4 cos 2theta wrong;
%   p. 49: « B_3 B_4 » in L_4; p. 56: det = (-1)^{n+1} rho;
%   p. 73: 3°) reads L_i not-in H_{i+1};
%   p. 81: Z_i ⊂ Z'_i iff h_{i+1} IN H_i (page: not-in); p. 82 c): irreducibility
%     on C needs rho != 0 as well, and the characteristic condition goes with
%     n EVEN; p. 80 (c): the two statements disagree, not settled;
%   p. 85: s_{i+1} NOT in X_i; p. 96 (14): E_i = k(s_0..s_i);
%   p. 101: the unit attached to <a_{i+1}, a'_i> = beta_i is not given and not
%     needed; the proof is redone;
%   p. 107-109: the recursion w(f_{V_{n+1}}) = (1+...) w(f_{V_n}) drops the
%     scaling of the sub-form by beta_0; the general product formula is wrong
%     as written (real counterexample n = 2, beta_0 = -1, beta_1 = 1); the
%     correct diagonalisation and formula are given; p. 110's second n = 2
%     computation is right and agrees with it; NB w_{n+1} = 0 is right;
%   p. 111: eps, eta_0..eta_2 are independent but not a basis (d(2) is
%     missing); the ten products are not independent (eta^2 = eps eta).
%
% Modern names given and footnoted as ours: abstract regular polytope and
% Schläfli symbol, string Coxeter group, Tits' geometric representation /
% Vinberg's linear Coxeter groups, Cartan and Gram matrices, continuants and
% the Desnanot-Jacobi identity, Schläfli's determinant criterion, circumsphere
% and midsphere, unitary reflection groups, the group schemes of
% Sekiguchi-Oort-Suwa, Stiefel-Whitney classes of quadratic forms (Delzant),
% Hasse-Witt and Clifford invariants, Brauer class of a conic bundle, Steinberg
% relation.
%
% Announced and never established: the Proposition « ?? » of p. 33 and its
% sequel on p. 34; the Théorème of p. 38 (no proof); the equivalence of p. 76
% (no proof); the independence question of pp. 111-112 and the Problème of
% p. 112; the end of p. 82 (« Ces sous- »), of p. 25, of p. 70.
\documentclass[11pt,a4paper]{article}
\input{../preamble/grothendieck}

\folder{72}
\pages{1}{112}
\dating{[à partir de 1978]}
\watermark{Édition de démonstration\\interprétation personnelle de l'œuvre}
\foldertitle{Polyèdres réguliers et hyperpolyèdres réguliers (Calculs en dimension quelconque) : notes manuscrites (s.d.) — lecture modernisée du dossier entier}

\begin{document}

\section*{Résumé}

\begin{resume}
Un kaléidoscope fabrique, avec deux ou trois miroirs et un seul objet, une
figure parfaitement symétrique. Les polyèdres réguliers --- le cube,
l'icosaèdre, et leurs analogues en dimension quelconque --- s'obtiennent de la
même façon : on part d'un sommet, on applique des réflexions dans un petit nombre
de miroirs, et les images successives du sommet sont tous les sommets. Le
dossier traite cette fabrication comme un problème d'algèbre linéaire, et le
traite en grand : les miroirs ne vivent plus dans l'espace usuel mais au-dessus
d'une base quelconque, comme si chaque coefficient était une fonction et non un
nombre, de sorte qu'on voit d'un coup tous les polyèdres réguliers d'un même
type et la façon dont ils dégénèrent les uns dans les autres.

La donnée de départ est minimale : $n + 1$ réflexions $\tau_0, \ldots,
\tau_n$, dont deux commutent dès qu'elles ne sont pas voisines dans la liste,
et un point de départ. Le résultat central est qu'une telle configuration,
pourvu qu'elle ne soit pas dégénérée, est décrite exactement par une poignée de
nombres : pour chaque paire de miroirs voisins, un nombre $\beta_i$ qui mesure
leur angle --- pour un cube, $2$ et $1$, parce que les faces sont des carrés
et que trois faces se rencontrent en chaque sommet. Tout le reste en découle
par des formules explicites : la forme quadratique invariante, qui est la
distance euclidienne ; le centre du polyèdre ; la sphère circonscrite et son
rayon ; le polyèdre dual ; le déterminant de la forme. Ces formules sont des
polynômes d'un type que les algébristes du dix-huitième
siècle connaissaient déjà, les \emph{continuants}, déterminants des matrices
qui n'ont de coefficients que sur trois diagonales.

Un de ces nombres, $\rho$, joue un rôle à part : il dit si le centre est à
distance finie. Quand $\rho$ s'annule, le centre part à l'infini, le polyèdre
devient un pavage du plan ou de l'espace --- le carrelage en carrés, le
pavage en hexagones ---, et le groupe des symétries qui fixent la figure
« à l'infini » passe d'un groupe de rotations-homothéties à un groupe de
translations. Le dossier suit cette dégénérescence pas à pas, ainsi que celles
qui n'ont lieu qu'en caractéristique $2$, où la symétrie par rapport au centre
--- l'« antipodisme » --- peut devenir l'identité.

Les dernières pages passent de la géométrie à la topologie : la forme
invariante a des invariants cohomologiques, ses classes de Stiefel-Whitney, et
il les calcule, avec une erreur que ses propres pages permettent de corriger.
Pour les polyèdres de l'espace, la question devient : la conique invariante
est-elle une vraie droite projective tordue, et la réponse est une classe de
Brauer.

Le dossier n'est pas rédigé ; il est fait de calculs, de recommencements et
de quelques énoncés nets --- un théorème sur les sous-espaces invariants, un
théorème de structure pour les familles de pseudo-réflexions. Les noms
modernes à chercher sont ceux des polytopes réguliers abstraits, des groupes de
Coxeter « en chaîne » et de leurs représentations par réflexions, des
continuants, des sphères circonscrites, et des invariants des formes
quadratiques.

\keywords{abstract regular polytope, Schläfli symbol, string Coxeter group, reflection representation, Tits representation, linear Coxeter group, pseudo-reflection, generalized Cartan matrix, Gram matrix, continuant, Desnanot-Jacobi identity, Schläfli determinant, circumscribed sphere, midsphere, dual polytope, antipodal map, unitary reflection group, group scheme deforming Gm to Ga, torsor of flags, invariant subspace, half-discriminant, Stiefel-Whitney class, Hasse-Witt invariant, Clifford invariant, Brauer class, conic bundle, Steinberg relation}
\end{resume}

\pagerange{2}{112}
\section*{Le fil du dossier, et les conventions}

\subsection*{Les feuillets}

Le dossier compte cent douze pages, sans pagination de sa main ; il numérote
seulement ses formules et ses conditions, par séries. Onze feuillets de garde
portent un titre de sa main en haut à droite et rien d'autre (pages 1, 26, 31,
36, 43, 55, 71, 78, 83, 95, 106) ; leurs titres sont repris ici comme titres de
sections. La page 21 est une garde blanche. Les pages 59, 61, 63, 65, 67 et 69
sont des pièces dactylographiées de l'université de Montpellier --- des avis de
soutenance de juin et juillet 1978, un recueil de décisions du Conseil de
l'Université de mai 1978 ---, sans rien de sa main ; elles s'accordent avec la
datation « à partir de 1978 » et n'en disent pas plus. La page 30 n'a pas de
marque dans la transcription, et l'en-tête de celle-ci n'en rend pas compte ;
cette lecture ne peut pas dire ce qu'elle porte.

\subsection*{Les stations}

\begin{itemize}
  \item \textbf{Pages 2 à 20 : le dictionnaire.} Une chaîne de catégories
    équivalentes décrivant les réalisations d'un polytope régulier ; les
    noyaux des actions et l'antipodisme (premier état, pages 5 à 9) ; la
    dualité (page 10) ; la torsion par un torseur (pages 11 et 12) ; un
    nouveau départ, avec la forme quadratique invariante calculée (pages 13 à
    20).
  \item \textbf{Pages 22 à 42 : l'antipodisme et le groupe $\mathcal{U}_\rho$.}
    L'antipodisme au second état (pages 22 à 25) ; une variante hermitienne
    (pages 27 à 29) ; la donnée « pseudo-sphérique » et la reconstruction des
    réflexions à partir des rotations (pages 32 à 35) ; le schéma en groupes
    $\mathcal{U}_\rho$ (pages 37 à 42).
  \item \textbf{Pages 44 à 60 : discriminants.} La matrice de la forme, les
    continuants $L_k$, le centre, $\rho$ et $\sigma$, la lissité, le plan
    antipodal ; des « calculs antérieurs » (pages 56 à 60).
  \item \textbf{Pages 62 à 77 : le dual et les calculs centrés.} L'identité
    $\sigma\sigma^\vee = \beta + \gamma\rho$, les sommets du dual ; puis, le
    centre étant supposé à distance finie, la sphère circonscrite, la sphère
    médiane, le discriminant.
  \item \textbf{Pages 79 à 82 : les sous-espaces invariants}, un théorème
    complet sur un corps.
  \item \textbf{Pages 84 à 105 : les fondements.} Réalisations d'une $n$-carte
    régulière, conditions de non-dégénérescence, torseur des drapeaux,
    coordonnées universelles, un lemme de relèvement, un théorème de structure
    pour les suites de pseudo-réflexions et l'existence du point de départ.
  \item \textbf{Pages 107 à 112 : classes de Stiefel-Whitney et classe de
    Brauer} de la forme invariante.
\end{itemize}

L'ordre des feuillets est celui des archivistes. Logiquement, les pages 84 à
105 viennent d'abord : elles posent ce que les autres calculent. On les cite
comme telles.

\subsection*{Conventions, valables pour tout le dossier}

\emph{La donnée.} $S$ est un schéma de base et $n \geqslant 1$ un entier.
$\mathfrak{G}_n$ est le groupe engendré par $\tau_0, \ldots, \tau_n$ soumis aux
seules relations $\tau_i\tau_j = \tau_j\tau_i$ pour $i, j$ non consécutifs ;
ajouter $\tau_i^2 = 1$ donne le groupe de Coxeter « en chaîne » dont tous les
ordres $\tau_i\tau_{i+1}$ sont libres, groupe d'automorphismes du polytope
régulier universel de rang $n+1$\footnote{Ce sont les noms d'aujourd'hui :
polytopes réguliers abstraits, groupes de Coxeter linéaires (« string Coxeter
groups ») de symbole de Schläfli $\{\infty, \ldots, \infty\}$ ; voir P.~McMullen
et E.~Schulte, \emph{Abstract Regular Polytopes} (2002). Les feuillets disent
« $n$-hyperpolyèdre combinatoire épinglé universel », « $n$-carte régulière »,
« quasi-polyèdre » ; le « groupe cartographique » engendré par $\sigma_0,
\ldots, \sigma_n$ de la page 84 est celui des dossiers 76 et 88.}. L'indice
$n$ est celui des $n$-cartes du dossier 76 : la carte est de dimension $n$, le
polytope convexe qui la borde de dimension $n+1$ ; pour les polyèdres de
l'espace usuel, $n = 2$.

$\mathcal{E}$ (on écrit $E$ quand rien n'est ambigu) est un fibré vectoriel de
rang $n+2$ sur $S$, $\pi \in E^\vee$ une forme linéaire partout non nulle,
$V = \operatorname{Ker} \pi$ de rang $n+1$. $X = \mathbb{P}(E)$ est le fibré
projectif dont les points sont les droites de $E$\footnote{Il écrit tantôt
$\mathbb{P}(E)$, tantôt $\check{\mathbb{P}}(E)$, selon la convention ; on fixe
celle-ci une fois pour toutes. « Rang » d'un fibré projectif, chez lui, est la
dimension relative : $X$ est « de rang $n+1$ ».}, $C = \mathbb{P}(V) \subset
X$ l'hyperplan à l'infini, $X \smallsetminus C$ le fibré affine des
$x$ tels que $\pi(x) = 1$. Une pseudo-réflexion de $E$ s'écrit
\[
\tau = \mathrm{id} + a' \otimes a, \qquad \tau(x) = x + \langle x, a' \rangle\, a ,
\]
avec $a \in E$, $a' \in E^\vee$ ; son \emph{centre} est $h = $ la droite de
$a$, son \emph{cocentre} l'hyperplan $H = \operatorname{Ker} a'$, sa valeur
propre $\lambda = 1 + \langle a, a' \rangle$ (inversible). C'est une réflexion
quand $\lambda = -1$.

\emph{Les vecteurs.} Le point de départ est $s_0 \in E$, $\pi(s_0) = 1$ ; on
pose $s_{i+1} = \tau_i s_i$, $a_i = s_{i+1} - s_i$ ($0 \leqslant i \leqslant
n$), $a_{-1} = s_0$ ; on normalise $\tau_i = \mathrm{id} + a'_i \otimes a_i$.
Les accouplements sont
\[
\langle a_{i-1}, a'_i \rangle = 1, \qquad
\langle a_i, a'_i \rangle = \lambda_i - 1 = \mu_i, \qquad
\langle a_{i+1}, a'_i \rangle = \beta_i, \qquad
\langle a_j, a'_i \rangle = 0 \ \text{sinon},
\]
et l'on pose $\alpha_i = \beta_i - 2$. Dans le cas des réflexions ($\lambda_i
= -1$), $\alpha_i$ est la trace de la « rotation » $\tau_i\tau_{i+1}$ sur le
plan qu'elle fait tourner, et, dans le cas euclidien, si $\tau_i\tau_{i+1}$
est d'ordre $p$,
\[
\beta_i = 4\cos^2\frac{\pi}{p}, \qquad \alpha_i = 2\cos\frac{2\pi}{p}.
\]
Pour le cube, de symbole de Schläfli $\{4, 3\}$, $(\beta_0, \beta_1) = (2,
1)$\footnote{Le lien avec le symbole de Schläfli est de cette lecture ; la
page 33 a $\beta_i = 4\cos^2\theta_i$ et appelle $\alpha_i$
« l'invariant de $\tau_i\tau_{i+1}$ (rotation) ». La matrice $(\langle a_j,
a'_i \rangle)$ est, au signe des $a'_i$ près, une matrice de Cartan
généralisée de type « chaîne », et le produit de ses deux coefficients non
diagonaux vaut $\beta_i$ : c'est la représentation géométrique de Tits
(Bourbaki, \emph{Groupes et algèbres de Lie}, ch.~V, \S\,4, 1968), sous la forme
plus libre des groupes de Coxeter linéaires de Vinberg (1971). Il pouvait
connaître Bourbaki et le livre de Coxeter, \emph{Regular Polytopes} ; aucune
page ne les cite.}.

\emph{Les continuants.} Pour $k \geqslant 0$, $L_k(B_1, \ldots, B_{k-1})$ est
le déterminant d'ordre $k$ de la matrice tridiagonale de diagonale $2$ dont
les coefficients sous- et sur-diagonaux ont pour produits $B_1, \ldots,
B_{k-1}$ :
\[
L_0 = 1, \quad L_1 = 2, \quad L_2(B_1) = 4 - B_1, \quad
L_k(B_1, \ldots, B_{k-1}) = 2L_{k-1}(B_2, \ldots) - B_1 L_{k-2}(B_3, \ldots).
\]
$L_k$ est symétrique ($L_k(B_1, \ldots, B_{k-1}) = L_k(B_{k-1}, \ldots, B_1)$),
pair pour $k$ impair, et l'on note $L'_k = L_k$ si $k$ est pair, $L_k/2$ si
$k$ est impair. On pose $\rho_{ij} = L_{j-i+2}(\beta_i, \ldots, \beta_j)$ (avec
$\rho_{i,i-1} = 2$, $\rho_{i,i-2} = 1$) et, pour tout le dossier,
\[
\rho = L_{n+1}(\beta_0, \ldots, \beta_{n-1}) = \rho_{0,n-1}, \qquad
\sigma = L_n(\beta_1, \ldots, \beta_{n-1}) = \rho_{1,n-1},
\]
$\tilde\rho = L'_{n+1}(\ldots)$, $\tilde\sigma = L'_n(\ldots)$ : l'un des deux
seulement est divisé par $2$, de sorte que $\rho\sigma = 2\tilde\rho\tilde\sigma$
toujours, $\rho = 2\tilde\rho$, $\sigma = \tilde\sigma$ si $n$ est pair, $\rho
= \tilde\rho$, $\sigma = 2\tilde\sigma$ si $n$ est impair. Enfin $\beta =
\beta_0 \cdots \beta_{n-1}$, $B = \beta_0^n \beta_1^{n-1} \cdots \beta_{n-1}$,
$\sigma^\vee = L_n(\beta_0, \ldots, \beta_{n-2}) = \rho_{0,n-2}$ et $\gamma =
L_{n-1}(\beta_1, \ldots, \beta_{n-2}) = \rho_{1,n-2}$\footnote{Les lettres
bougent d'une page à l'autre. Les « moitiés » $\tilde\rho, \tilde\sigma$ sont
ses $\rho', \sigma'$ aux pages 22 à 24, 52 à 54 et 74 à 77, ses $\tilde\rho,
\tilde\sigma$ aux pages 39 à 42, 62 à 66 et 107 ; mais aux pages 60 à 70 son
$\sigma'$ est $\rho_{0,n-2}$, qu'on note ici $\sigma^\vee$. On garde une lettre
par objet.}.

\emph{Le centre.} Le vecteur $\tilde c \in E$ est l'unique vecteur annulé par
tous les $a'_i$ dont la coordonnée sur $a_n$ vaut $1$ ; on montrera (pages 47
et 51) que $\pi(\tilde c) = \rho$. Le centre du polytope est le point $c$ de
$X$ qu'il définit ; il est à distance finie si et seulement si $\rho$ est
inversible, et alors $c = \tilde c/\rho$ dans l'espace affine ; on pose $c' =
2c$\footnote{Sa lettre pour ce vecteur varie : $\tilde\sigma$ (page 22),
$c'$, $\tilde c$ (pages 47 à 53), $\pi$ (pages 5 à 9, 38 à 40, 60 à 70, 107),
et une sorte de $\sigma$ bouclé (page 77). Réciproquement, la forme $\pi$ de
cette lecture est son $\pi$ aux pages 22 à 25 et son $\pi'$ ailleurs.}.

\emph{La forme.} $f$ est une forme quadratique, $\varphi(x, y) = f(x+y) -
f(x) - f(y)$ sa forme polaire, $\psi \colon E \to E^\vee$, $x \mapsto
\varphi(x, \cdot)$. On note $f_V$ la forme invariante sur $V$ normalisée par
$f_V(a_0) = 1$, $f_E$ son extension à $E$ normalisée par $f_E(s_0) = 0$, et
$c_i = f(a_i)$\footnote{Son indice « $\varepsilon$ » ($f^\varepsilon$,
$f_\varepsilon$, $\delta_\varepsilon$, pages 41, 42 et 52) est lu ici comme un
$\mathcal{E}$ cursif : c'est la forme sur $\mathcal{E}$, ce que confirment les
discriminants de la page 52.}.

\subsection*{Ce que seule une lecture d'ensemble peut dire}

Les pages 84 à 105 donnent la raison pour laquelle le reste du dossier peut
calculer comme il le fait. Le théorème des pages 96 à 100, appliqué dans les
« coordonnées universelles » de la page 91, montre qu'une réalisation
projective non dégénérée et épinglée est rigide et entièrement déterminée par
ses $2n+1$ invariants $(\lambda_0, \beta_0, \lambda_1, \ldots, \beta_{n-1},
\lambda_n)$, les $\lambda_i$ inversibles et les $\beta_i$ quelconques ---
c'est l'affirmation, sans démonstration, de la note de marge de la page 13.
Le groupe $\mathcal{U}_\rho$ des pages 37 à 39 et le groupe
$\underline{\mathrm{Aut}}(D_0, c + d)$ des pages 87 et 88 sont le même objet
vu deux fois : le second est le groupe des automorphismes du système des
$\tau_i$, le premier en est la description explicite, et la dégénérescence
de $\mathbb{G}_m$ en $\mathbb{G}_a$ y est le passage du centre à l'infini,
$\rho \to 0$. La reconnaissance du symbole de Schläfli dans les $\beta_i$, et
celle de $\rho$ comme un multiple du déterminant de Schläfli, sont de cette
lecture.

\subsection*{Ce que le dossier annonce sans l'établir}

La proposition marquée « ?? » de la page 33, et sa suite de la page 34 ; le
« Théorème » de la page 38, énoncé sans preuve ; l'équivalence de catégories de
la page 76, énoncée sans preuve ; l'indépendance des dix classes de la page 111
et le « Problème » de la page 112 ; la fin des pages 25, 70 et 82.

\pagerange{2}{4}
\section*{I. Le dictionnaire (pages 1 à 4)}

La garde porte : « $n$-polyèdres réguliers et représentations de
$\mathfrak{G}_n$ dans $\mathrm{O}(n+1)$ (dictionnaire) »\footnote{Le signe
de « $n+1$ » est d'une lecture incertaine ; à gauche, au crayon, « (23 p) »,
d'une main non établie.}. Le dictionnaire est une liste de catégories, toutes
équivalentes, qui décrivent la même chose --- les réalisations géométriques
régulières non dégénérées du polytope universel épinglé $\Pi$ --- en mettant
l'accent sur des structures différentes :
\begin{itemize}
  \item $C_{\mathrm{pr}}$, les réalisations projectives ; $C'_{\mathrm{pr}}$,
    les systèmes $(X, (\tau_i), X_*)$ d'un fibré projectif, de
    pseudo-réflexions strictes et d'un drapeau maximal $X_*$ tel que $\tau_i
    X_j = X_j$ pour $i \neq j$ et $\tau_i X_i \neq X_i$ ;
  \item $C_{\mathrm{aff}}, C'_{\mathrm{aff}}$, leurs variantes affines, et
    $C_{\mathrm{vect}}, C'_{\mathrm{vect}}$, leurs variantes vectorielles, où
    $X = \mathbb{P}(E)$ et le point de départ est rigidifié en un vecteur ;
  \item $C^\infty_{\mathrm{pr}}$ et $C^\infty_{\mathrm{vect}}$, les données
    « à l'infini » : un fibré projectif $C$ de dimension relative $n$ (ou un
    fibré vectoriel $V$ de rang $n+1$), une action de $\mathfrak{G}_n$ dont
    les $\tau_i$ sont des pseudo-réflexions, strictes pour $i \geqslant 1$,
    et un point $h_0$ (une droite $L_0$), avec les conditions : les centres
    engendrent, et $h_i \notin H_{i+1}$ sur chaque fibre ;
  \item $C^\infty_{\mathrm{pr.orth}}$ et $C^\infty_{\mathrm{vect.orth}}$, les
    mêmes avec une hyperquadrique invariante $Q$ (une forme quadratique $f$)
    et des réflexions orthogonales, le point de départ normalisé par $f(a_0)
    = 1$.
\end{itemize}
Il note lesquelles sont autoduales (les projectives et vectorielles) et
lesquelles ne le sont pas « sauf si $\rho$ inversible » (les données à
l'infini). La chaîne des équivalences s'écrit
\[
C_{\mathrm{pr}} \simeq C'_{\mathrm{pr}} \simeq C_{\mathrm{aff}} \simeq
C'_{\mathrm{aff}} \simeq C_{\mathrm{vect}} \simeq C'_{\mathrm{vect}}
\simeq C^\infty_{\mathrm{pr}} \simeq C^\infty_{\mathrm{vect}}
\simeq C^\infty_{\mathrm{pr.orth}} \simeq C^\infty_{\mathrm{vect.orth}},
\]
où le passage aux données à l'infini demande $n \geqslant 2$, et le passage
aux données orthogonales que les $\tau_i$ soient des réflexions ($\lambda_i
= -1$)\footnote{Sur la page, deux traits verticaux coupent la chaîne ; au
premier est accrochée la mention « $n \geqslant 2$ (explicités variante si $n =
1$ !) », au second « si de plus les $\lambda_i = -1$ i.e.\ les $\tau_i$ des
réflexions ». La raison de la première restriction est le lemme de relèvement
de la page 94 : voir la section IV.}.

Trois remarques l'accompagnent. Les catégories $C$ proviennent par
\emph{rigidification} de catégories $\widetilde{C}$ de représentations
projectives de $\mathfrak{G}_n$ « sans plus » : rigidifier, c'est se donner le
point de départ (ou une base de la droite $L$ commune aux isomorphismes
$L_0 \simeq L_1^\vee \simeq L_1 \simeq \cdots \simeq L_n$), ce qui tue le
groupe $\mathbb{G}_m$ des automorphismes. Dans la variante affine, $\tau_0$
n'est plus forcément une pseudo-réflexion affine, mais son prolongement
projectif l'est, son hyperplan pouvant être l'hyperplan à l'infini. Dans le cas
orthogonal, si $\tau_i$ et $\tau_{i+1}$ sont des réflexions orthogonales
strictes, la condition $h_i \notin H_{i+1}$ équivaut à ce que $\tau_i$ et
$\tau_{i+1}$ ne commutent sur aucune fibre, c'est-à-dire à
$(\tau_i\tau_{i+1})^2 \neq \mathrm{id}$ sur chaque fibre. Dans
$C^\infty_{\mathrm{vect.orth}}$ la rigidification $f(a_0) = 1$ n'est
définie qu'au signe près : les sections possibles forment un torseur sous
$\mu_2$.

\pagerange{5}{9}
\section*{II. Noyaux et antipodisme, premier état (pages 5 à 9)}

Soit $G_0$ l'image de $\mathfrak{G}_n$ dans le groupe des automorphismes de
$(E, f, \pi, \tilde c)$. Il cherche les noyaux des flèches
\[
\operatorname{Aut}(E, f, \pi, \tilde c) \longrightarrow
\operatorname{Aut}(X), \quad \operatorname{GL}(V), \quad
\operatorname{PGL}(V) = \operatorname{Aut}(C).
\]

\emph{Vers $\operatorname{Aut}(X)$.} Le noyau de $\operatorname{Aut}(E, f)
\to \operatorname{Aut}(X)$ est formé des homothéties qui conservent $f$, soit
$\mu_2$ ; une homothétie qui fixe en plus $\tilde c$ (ou $\pi$) est
l'identité. Donc $G_0$ opère fidèlement sur $X$.

\emph{Vers $\operatorname{GL}(V)$.} Un automorphisme qui est l'identité sur
$V$ s'écrit $u = \mathrm{id} + \pi \otimes a$, $u(x) = x + \pi(x)\, a$. Il
conserve $\pi$ si et seulement si $\pi(a) = 0$, et fixe $\tilde c$ si et
seulement si $\rho a = 0$ ; si $\rho$ n'est pas diviseur de zéro, $u =
\mathrm{id}$. La Proposition de la page 6 traite le reste : sur un anneau
intègre, la seule exception est celle où l'antipodisme (section VI) fixe
$\tilde c$. Or l'antipodisme $u = \mathrm{id} - 2\pi \otimes c$ est
l'identité sur $V$ et envoie $\tilde c$ sur $-\tilde c$ ; il fixe donc
$\tilde c$ exactement en caractéristique $2$, et il n'y est pas l'identité
exactement quand $n$ est pair --- ce sont les conditions « $n$ pair, car.\ $2$,
$\tilde\rho$ inversible » de l'énoncé\footnote{L'énoncé de la page 6 s'écrit
« Sur un anneau intègre, si $\beta$ ou $\rho \neq 0$, alors le seul $u$
\ldots\ induisant l'identité dans $V$ est l'identité, sauf dans le cas $\beta
\neq 0$, $n$ pair, car.\ 2 ($\rho = 2\rho' = 0$), $\rho'$ inv. » La lecture
ci-dessus en donne la raison ; les calculs des pages 5 à 8 sont menés dans des
notations provisoires ($\beta$, $\sigma$, $\sigma'$, $\psi'$) qui ne sont pas
celles des pages suivantes, et en partie barrés (la moitié supérieure de la
page 8 l'est, et l'alinéa entre crochets des pages 7 à 9 est un calcul
abandonné en coordonnées).}.

\emph{Vers $\operatorname{PGL}(V)$.} Un automorphisme qui induit l'identité
sur $C = \mathbb{P}(V)$ induit sur $V$ une homothétie $\varepsilon$ qui
conserve $f|_V$, donc $\varepsilon \in \mu_2$, et $\varepsilon u$ est du cas
précédent. Sur un corps, $\varepsilon = 1$ ramène au cas déjà traité ;
$\varepsilon = -1 \neq 1$ (caractéristique $\neq 2$) force $\rho a = 2\tilde
c$, donc $\rho$ inversible et $\varepsilon u = -\mathrm{id} + \pi \otimes
(2/\rho)\tilde c = -\mathrm{id} + \pi \otimes c'$ : c'est l'antipodisme, dont
on vérifie qu'il conserve $f$, $\pi$ et la droite de $\tilde c$. D'où la
Proposition des pages 8 et 9 : \emph{sur un corps, sous l'hypothèse « $\beta$ ou
$\rho$ inversible » de la page --- où $\beta$ mesure la non-dégénérescence de
$f|_V$ ---, le seul automorphisme non trivial de $(E, f, \pi, \pm\tilde c)$ qui
induise l'identité sur $C$ est l'antipodisme}, défini dès que $\tilde\rho$ est inversible, par
\[
u = \mathrm{id} - \frac{1}{\tilde\rho}\, \pi \otimes \tilde c \quad (n \text{
pair}), \qquad
u = \mathrm{id} - \frac{2}{\rho}\, \pi \otimes \tilde c \quad (n \text{
impair}),
\]
le second étant l'identité en caractéristique $2$\footnote{Sur la page, le
premier cas porte $\pi' \otimes \pi$ dans ses notations, où $\pi'$ est notre
forme $\pi$ et $\pi$ notre vecteur $\tilde c$. Les deux formules sont la même,
$u = \mathrm{id} - \pi \otimes c'$, puisque $c' = \tilde c/\tilde\rho$ pour $n$
pair et $2\tilde c/\rho$ pour $n$ impair.}. Projectivement, $u$ est la symétrie
$x \mapsto 2c - x$ par rapport au centre.

\pagerange{10}{12}
\section*{III. Dualité et torsion (pages 10 à 12)}

\emph{Dualité} (page 10). Si $\rho$ est inversible, l'accouplement naturel
entre $V$ et le $V'$ du système dual est parfait, donc $C$ et $C'$ sont en
dualité. Si de plus $f|_V$ est non dégénérée, $\psi_V \colon V \simeq V'$
donne un isomorphisme $C \simeq C'$ qui commute à $\mathfrak{G}_n$ et
transforme $Q_C$ en $Q_{C'}$. Le passage au polytope dual s'exprime alors, dans
$C^\infty_{\mathrm{pr.orth}}$, par la conservation de $(C, Q_C)$ et le
renversement $\tau_i \mapsto \tau_{n-i}$, c'est-à-dire en composant
$\varphi^C$ avec l'automorphisme $\theta$ de $\mathfrak{G}_n$ qui échange
$\tau_i$ et $\tau_{n-i}$ ; il n'y a plus besoin de point de départ rigidifié
puisque $\beta_0$ et $\beta_{n-1}$ sont inversibles\footnote{La page écrit
« $\beta_0, \beta_n$ » ; dans les conventions de cette lecture le dernier
paramètre est $\beta_{n-1}$.}. C'est la dualité polaire des polytopes convexes
(dossier 87), lue sur les réflexions.

\emph{Torsion} (pages 11 et 12). Soit maintenant $\Pi$ un polytope régulier
combinatoire de dimension $n$, $G$ son groupe d'automorphismes et $R_\Pi$ le
torseur de ses repères (ses drapeaux) ; le commutant de $G$ dans les
permutations de $R_\Pi$ est un quotient $\mathfrak{G}_\Pi$ de $\mathfrak{G}_n$.
Chaque repère $r$ donne un objet de $C_{\mathrm{pr}}$, et ces objets sont reliés
par un système transitif d'isomorphismes ; donc les réalisations de $\Pi$
forment la sous-catégorie pleine de $C_{\mathrm{pr}}$ des objets dont la
représentation se factorise par $\mathfrak{G}_\Pi$. On retrouve la réalisation
de $\Pi$ par le produit contracté
\[
R_\Pi \wedge^{\mathfrak{G}_\Pi} X ,
\]
$\mathfrak{G}_\Pi$ opérant à droite sur $R_\Pi$\footnote{Le produit contracté
est celui du dossier 75. La page parle de « tordre » ; le mot « commutant »
est lu avec doute.}. Autre point de vue : on se donne seulement le quotient
$\mathfrak{G}_0$ de $\mathfrak{G}_n$ correspondant à un polytope régulier
épinglé $\Pi_0$ ; un polytope régulier combinatoire « de type $\mathfrak{G}_0$ »
sur $S$ est un tordu de $\Pi_0$ par un torseur sous $\mathfrak{G}_0$, et la
catégorie de ses réalisations projectives est équivalente au produit de la
catégorie des réalisations de type $\mathfrak{G}_0$ par celle des
$\mathfrak{G}_0$-torseurs\footnote{La page dit « produit de $C_{\mathrm{pr}}$
(ou d'une de ses variantes) et de la cat.\ des $\mathfrak{G}_0$-torseurs » ;
c'est la sous-catégorie pleine de type $\mathfrak{G}_0$ de la page 11 qui
convient. La lettre $\mathfrak{G}_S$ de la page 12 est lue avec doute.}.

\pagerange{13}{18}
\section*{IV. Nouveau départ : les catégories $C_0$ à $C'''_0$ et la forme invariante (pages 13 à 18)}

\emph{$C_0$} (page 13) est le groupoïde des réalisations projectives non
dégénérées : triples $(X, (\tau_i)_{0 \leqslant i \leqslant n}, s_0)$, $X$ fibré
projectif de dimension relative $n+1$, $\tau_i$ pseudo-réflexions strictes de
centres $h_i$ et cocentres $H_i$, $s_0$ une section, avec
\begin{enumerate}
  \item[1)] les $h_i$ linéairement indépendants sur chaque fibre (ils
    engendrent donc un hyperplan $C$) ;
  \item[1')] les $H_i$ indépendants sur chaque fibre (leur intersection est
    donc une section $c$) ;
  \item[2)] $\tau_i\tau_j = \tau_j\tau_i$ pour $i, j$ non consécutifs ;
  \item[3)] $h_i \notin H_{i+1}$ sur chaque fibre ($0 \leqslant i \leqslant
    n-1$) ;
  \item[4)] $s_0 \in H_i$ pour $1 \leqslant i \leqslant n$ ;
  \item[5)] $s_0 \notin H_0$ sur chaque fibre.
\end{enumerate}
Il note que 1') suit des autres conditions en présence de $s_0$, mais qu'il
faut l'imposer sans lui. Et, en marge : \emph{$C_0$ est une catégorie rigide,
dont les classes d'isomorphie correspondent aux suites $(\lambda_0, \alpha_0,
\lambda_1, \ldots, \alpha_{n-1}, \lambda_n)$ de sections de
$\mathcal{O}_S$}\footnote{La page écrit l'exposant « $2n-1$ » (lecture
incertaine) ; la liste qu'il compte a $2n+1$ termes. Les $\lambda_i$, valeurs
propres de pseudo-réflexions, sont des unités. La démonstration de cette
assertion n'est pas sur la page ; elle découle du théorème des pages 96 à 100
lu dans les coordonnées universelles de la page 91 (section XV).}.

\emph{$C'_0$} (pages 13 et 14) est la catégorie des triples $(V,
\varphi^V, a_0)$ : $V$ fibré vectoriel de rang $n+1$, $\varphi^V \colon
\mathfrak{G}_n \to \operatorname{GL}(V)$, $a_0$ une section partout non nulle,
avec : 1) $\tau_i^V$ est une pseudo-réflexion, stricte pour $i \geqslant 1$,
de centre $\mathcal{O}_S a_0$ pour $i = 0$ ; d'où des droites $L_i \subset V$,
$L'_i \subset V^\vee$ ; 2) $V = \bigoplus L_i$ ; 4) $L_i \not\subset
H_{i+1}$ sur chaque fibre, c'est-à-dire $L_i \otimes L'_{i+1} \simeq
\mathcal{O}_S$ par l'accouplement. De 1), $L'_i \otimes L_i \simeq \mathcal{O}_S$
($\tau_i - \mathrm{id} \mapsto 1$), donc $L'_i \simeq L_i^\vee$, et de proche
en proche
\[
L_0 \simeq L_1 \simeq \cdots \simeq L_n =: L, \qquad V \simeq L \otimes
\mathcal{O}_S^{n+1},
\]
de sorte que $V$ a une base canonique $(a_0, \ldots, a_n)$ déterminée par
$a_0$, avec $\tau_i = \mathrm{id} + a_i^\vee \otimes a_i$ et $\langle a_i,
a_{i+1}^\vee \rangle = 1$. La commutation des $\tau_i$ non consécutifs se
traduit par $L_i \subset \operatorname{Ker} a_j^\vee$ pour $i, j$ non
consécutifs, d'où, dans la base duale $(a_i^*)$,
\[
a_0^\vee = \mu_0\, a_0^* + \beta_0\, a_1^*, \qquad
a_i^\vee = a_{i-1}^* + \mu_i\, a_i^* + \beta_i\, a_{i+1}^*, \qquad
a_n^\vee = a_{n-1}^* + \mu_n\, a_n^*
\]
\footnote{Le coefficient de $a_{i-1}^*$ dans la ligne médiane est un signe
douteux sur la page ; c'est $1$, par $\langle a_{i-1}, a_i^\vee \rangle = 1$, et
la dernière ligne l'écrit sans coefficient.}. En marge : $\tau_0$ est stricte
si et seulement si $\mu_0\mathcal{O}_S + \beta_0\mathcal{O}_S = \mathcal{O}_S$.
La \emph{Scholie} de la page 15 affirme que le foncteur évident $C_0 \to C'_0$
est une équivalence.

\emph{$C''_0$} (pages 15 et 16) est la catégorie des triples $(C, \varphi^C,
h_0)$ à l'infini : $C$ de dimension relative $n$, les $\tau_i^C$
pseudo-réflexions, strictes pour $i \geqslant 1$, de centre $h_0$ pour $i =
0$, les $h_i$ indépendants, $h_i \notin H_{i+1}$. Le foncteur $C'_0 \to C''_0$,
$V \mapsto \mathbb{P}(V)$, est une équivalence \emph{si $n \geqslant 2$} :
localement $C = \mathbb{P}(V)$, et chaque pseudo-réflexion de $\mathbb{P}(V)$ se
relève de façon unique en une pseudo-réflexion de $V$ dès que $\operatorname{rg}
V \geqslant 3$ ; $V$ n'est déterminé qu'à un $\otimes L$ près, et la
rigidification lève l'indétermination. D'où la \emph{Rescholie} : pour $n
\geqslant 2$, $C_0 \simeq C'_0 \simeq C''_0$\footnote{La page écrit « si $n
\geqslant 2$ (mais non si $n = 2$) », l'un des deux portant peut-être un $1$.
La restriction juste est $n \geqslant 2$ : le relèvement unique est le lemme
de la page 94 (rang $\geqslant 3$), et il échoue sur la droite projective, où
une involution à deux points fixes se relève en deux pseudo-réflexions
différentes, selon le point fixe qu'on déclare « hyperplan ».}.

\emph{La forme invariante} (pages 16 et 17). Les formes quadratiques invariantes
sur $V$ partout non nulles correspondent aux hyperquadriques invariantes $Q_C$
de $C$. Pour qu'une forme $f$ soit invariante par la pseudo-réflexion $\tau =
\mathrm{id} + a' \otimes a$, il faut et il suffit (sur une base intègre, $a'
\neq 0$) que
\[
f(a)\, a' + \psi(a) = 0 ,
\]
car $f(\tau x) - f(x) = \langle x, a' \rangle\bigl(\varphi(x, a) + \langle x, a'
\rangle f(a)\bigr)$. Appliqué aux $\tau_i$, cela donne $\varphi(a_j, a_i) =
-c_i \langle a_j, a'_i \rangle$, donc $\varphi(a_i, a_j) = 0$ pour $i, j$ non
consécutifs, $\varphi(a_i, a_i) = -c_i\mu_i$, et, en écrivant la symétrie de
$\varphi$ sur $(a_i, a_{i+1})$,
\[
c_{i+1} = \beta_i c_i , \qquad \varphi(a_i, a_{i+1}) = -c_{i+1}.
\]
Comme $\varphi(a_i, a_i) = 2c_i$, la condition $c_i(\lambda_i + 1) = 0$ force
$\lambda_i = -1$ dès qu'on veut $f(a_i)$ inversible : il faut des réflexions,
et dans le groupe on peut alors ajouter les relations $\tau_i^2 = 1$. Sous
$\lambda_i = -1$, les formes invariantes forment un module libre de rang $1$,
engendré par
\[
f_V\Bigl(\sum_{i=0}^n x_i a_i\Bigr) = \sum_{i=0}^{n} (\beta_0 \cdots
\beta_{i-1})\, x_i^2 \;-\; \sum_{i=0}^{n-1} (\beta_0 \cdots \beta_i)\, x_i x_{i+1}
\]
\footnote{La page écrit les termes rectangles avec le signe $+$ ; ses propres
relations $\varphi(a_i, a_{i+1}) = -c_{i+1}$ donnent $-$, comme les pages 44 et
60, qui l'écrivent ainsi. On vérifie directement que cette $f_V$ satisfait
$f(a_i)a'_i + \psi(a_i) = 0$ pour tout $i$, sans aucune hypothèse
d'inversibilité.} : $f_V(a_i) = c_i = \beta_0 \cdots \beta_{i-1}$. Sur un corps,
il y a donc une seule hyperquadrique invariante $Q_C$, et $\mathfrak{G}_n$
opère sur $(C, Q_C)$. Pour le cube, $f_V$ est le carré de la longueur
euclidienne quand l'arête $a_0$ est de longueur $1$ : $f(a_1) = 2$ est le carré
de la diagonale d'une face.

\emph{$C'''_0$} (page 18) est la catégorie des quadruples $(C, Q_C, \varphi^C,
h_0)$, avec les conditions a) à c) de $C''_0$ et $\varphi^C$ à valeurs dans
$\operatorname{Aut}(C, Q_C)$. Il note que $\tau_0^C$ est une réflexion
orthogonale stricte si et seulement si, en chaque point, $\beta_0 \neq 0$ ou la
caractéristique résiduelle est $\neq 2$. Si les $\tau_i$ sont des réflexions
orthogonales strictes (les $\beta_i$ inversibles, ce qui est le cas si $Q_C$
est lisse), la condition c) équivaut à

c') $(\tau_i\tau_{i+1})^2 \neq \mathrm{id}$ sur chaque fibre, c'est-à-dire
que $\tau_i$ ne commute à $\tau_{i+1}$ sur aucune fibre\footnote{La page écrit
« $(\tau_i\tau_{i+1})^2 = \mathrm{id}$ dans chaque fibre, i.e.\ $\tau_i$ ne
commute pas à $\tau_{i+1}$ » : le signe est un lapsus, que la paraphrase et la
page 3 (« $\rho_i^2 \neq \mathrm{id}$ en tte fibre ») corrigent.}.

Et la condition b) s'énoncerait :

b') pour tout $s \in S$, la représentation induite sur $C_s$ est
irréductible.

Il ajoute : « suffisance claire, pour la nécessité il faut travailler un peu »,
avec un point d'interrogation en marge. De fait b') entraîne b) aussitôt, le
sous-espace engendré par les $h_i$ étant invariant ; mais la réciproque est
fausse en général. Par le théorème des pages 79 à 82 (section XIV), la
représentation sur $C_s$ est irréductible si et seulement si tous les
$\beta_i(s)$ et $\rho(s)$ sont non nuls ; sous les hypothèses de la remarque
($\beta_i$ inversibles), b) et b') sont donc équivalentes exactement quand le
centre n'est à l'infini sur aucune fibre\footnote{Le rapprochement est de
cette lecture ; le point d'interrogation de la page est justifié.}.

\pagerange{19}{20}
\section*{V. Le cas $n$ pair, et l'exemple des polyèdres de l'espace (pages 19 et 20)}

Pour $n$ pair, $V$ est de rang impair, donc, $f_V$ étant lisse, $\mathrm{O}(f_V)
= \mathrm{SO}(f_V) \times \mu_2$ et $\mathrm{SO}(f_V) \hookrightarrow
\operatorname{Aut}(C, Q_C)$. L'homomorphisme $\varphi^C \colon
\mathfrak{G}_n \to \operatorname{Aut}(C, Q_C)$ se relève alors à
$\mathrm{SO}(f_V)$ par
\[
\psi^V(\tau_i) = -\tau_i^V ,
\]
qui est de déterminant $+1$ puisque $\tau_i^V$ est une réflexion en rang
impair.

\emph{L'exemple $n = 2$} (polyèdres de l'espace). $f_V$ est lisse si et
seulement si $\beta_0$, $\beta_1$ et $\tilde\rho = 4 - \beta_0 - \beta_1 =
-(\alpha_0 + \alpha_1)$ sont inversibles\footnote{La page écrit « $\beta_0,
\beta_1$ et $\beta_0 + \beta_1 - 4 = \alpha_0 + \alpha_1$ inv. », ce qui est
la même condition.}. Alors $Y = Q_C$ est une conique lisse dans le plan
projectif $C$, un fibré en droites projectives sur $S$, et $\mathrm{SO}(f_V)
\simeq \underline{\operatorname{Aut}}(Y)$ : c'est l'isomorphisme exceptionnel
$\mathrm{SO}_3 \simeq \mathrm{PGL}_2$ sous sa forme tordue, objet du dossier 82.
Se donner $\varphi^C$ revient donc à se donner la conique $Y$ et l'action de
$\mathfrak{G}_2$ sur elle par $\tau^Y_0, \tau^Y_1, \tau^Y_2$, et les conditions
deviennent : 1) chaque $\tau_i^C$ est $\neq \mathrm{id}$ sur chaque fibre ;
2) $h_0, h_1, h_2$ sont indépendants sur chaque fibre, ce qu'il traduit par :
il n'y a sur aucune fibre géométrique de $Y$ de diviseur de degré $2$
invariant, c'est-à-dire pas de réduction de la structure à une droite affine
ou à un torseur sous $\mathbb{G}_m$ --- « c'est sans doute ce qui exclut le cas
diédral » ; 3) les rotations $\rho_f = \tau_0\tau_1$ et $\rho_s = \tau_1\tau_2$
vérifient $\rho_f^2 \neq \mathrm{id}$, $\rho_s^2 \neq
\mathrm{id}$ sur chaque fibre\footnote{Les indices $f$ (face) et $s$ (sommet)
sont lus avec doute. La traduction de 2) est la sienne et n'est pas
démontrée ; une direction au moins est claire par la polarité relative à la
conique, qui échange points et droites de $C$ : si les $h_i$ sont alignés sur
une droite $\ell$, celle-ci est invariante et $\ell \cap Y$ est un diviseur
invariant de degré $2$.}.

\pagerange{22}{25}
\section*{VI. L'antipodisme, second état (pages 22 à 25)}

La page 22 reprend la question des pages 5 à 9 avec les notations définitives :
quels automorphismes de $X$ sont l'identité sur $C$ et commutent aux
$\tau_i$ ? Ce sont les $u = \mathrm{id} + \pi \otimes a$ avec $\tau_i a = a$
pour tout $i$, donc $a$ multiple de $\tilde c$ :
\[
u_\lambda = \mathrm{id} + \lambda\, \pi \otimes \tilde c .
\]
Un calcul direct donne $f_E(u_\lambda x) = f_E(x) - \lambda\,\pi(x)^2(\sigma +
\lambda\tilde\sigma\tilde\rho)$, en utilisant $\psi(\tilde c) = -\sigma\pi$ et
$f_E(\tilde c) = -\tilde\sigma\tilde\rho$ (pages 47 et 52). Donc $u_\lambda$ est
orthogonal si et seulement si
\[
\lambda(\sigma + \lambda\tilde\sigma\tilde\rho) = 0, \quad \text{soit} \quad
\lambda\tilde\sigma(1 + \lambda\tilde\rho) = 0 \ (n \text{ pair}), \qquad
\lambda\tilde\sigma(2 + \lambda\tilde\rho) = 0 \ (n \text{ impair}).
\]
C'est la Proposition de la page 22, complétée page 23 : si $\tilde\sigma = 0$,
toutes les $u_\lambda$ sont orthogonales\footnote{La fin de l'alinéa 1) de la
Proposition se lit « sont \emph{nulles} », avec doute ; l'équation ne laisse que
la lecture donnée.} ; si $\tilde\sigma$ est inversible,
pour $n$ pair, il existe un tel $u$
non trivial sur chaque fibre si et seulement si $\tilde\rho$ est inversible, et
alors $\lambda = -1/\tilde\rho$ ; pour $n$ impair, il faut et il suffit que
$\tilde\rho$ et $2$ soient inversibles, et $\lambda = -2/\rho$. Dans les deux
cas $u = \mathrm{id} - \pi \otimes c'$ est \emph{l'antipodisme}\footnote{Les
équations (3), (4), (6), (7) de la page sont celles-ci, à la notation près
($\tilde\sigma$ y désigne le vecteur centre, et $\rho', \sigma'$ les moitiés).
Le cas c) de la page 23, sur la caractéristique $2$, est trop surchargé pour
être lu ; la condition « $u$ réflexion » de l'équation (5), $\lambda\rho = -2$,
est la même que celle de l'antipodisme.}.

\emph{Corollaire} (page 23). Supposons $\tilde\sigma\tilde\rho$ inversible,
c'est-à-dire $f_E(\tilde c)$ inversible, c'est-à-dire $c \notin Q_X$ sur
chaque fibre, et $2$ inversible si $n$ est impair. Alors il existe un unique
$\mathfrak{G}_n$-automorphisme de $X$, non trivial sur chaque fibre, au-dessus
de l'identité de $C$ et conservant $Q_X$ : l'antipodisme.

\emph{NB} (page 24). L'antipodisme est défini intrinsèquement dès que
$\tilde\rho$ est inversible ; il est l'identité exactement quand $n$ est impair
et la caractéristique $2$ ; il est donc non trivial sur chaque fibre si et
seulement si $n$ est pair ou $2$ inversible.

\emph{Reconstruction des bases} (pages 24 et 25). Inversement, partant de
$\varphi_G \colon \mathfrak{G}_n \to \operatorname{Aut}(X)$ avec les conditions
habituelles, on a $C \subset X$, l'extension $0 \to V \to E \to \mathcal{O}_S
\to 0$, des droites $L_i \subset V$, $L'_i \subset V^\vee$ avec $V =
\bigoplus L_i$ et $L_i \otimes L'_{i+1} \simeq \mathcal{O}_S$. La donnée de
$s_0$, section de $E$ au-dessus de $1$, fixe $a'_0$ par $\langle s_0, a'_0
\rangle = 1$, puis de proche en proche les bases $a_0, a'_1, a_1, \ldots,
a'_n, a_n$ par
\[
\tau_i = \mathrm{id} + a'_i \otimes a_i , \qquad \langle a_i, a'_{i+1}
\rangle = 1 .
\]
Réciproquement, $a'_0$ étant donné, le plan $F = \bigcap_{i \geqslant 1} H_i$
n'est contenu dans $H_0$ sur aucune fibre, $F \cap H_0$ est la droite de
$\tilde c$, et l'équation $\langle s, a'_0 \rangle = 1$ définit dans $F$ une
droite affine parallèle à celle-ci, lieu des points de départ possibles. La
page s'arrête là, sur la projection $\pi|_F$ ; la section XV en donnera la
suite (le torseur $D_0^{**}$).

\pagerange{27}{29}
\section*{VII. Forme hermitienne invariante pour un hyperpolyèdre régulier (pages 26 à 29)}

Le titre de la garde annonce une variante complexe : les $\tau_i$ sont des
pseudo-réflexions unitaires de valeurs propres $\lambda_i$ de module $1$ (ou,
sur une extension quadratique $k/k_0$, de norme $1$), et l'on cherche une forme
hermitienne $h$ invariante. Une pseudo-réflexion unitaire de vecteur $a$ et de
valeur propre $\lambda$ s'écrit $x \mapsto x + (\mu/h(a, a))\, h(x, a)\, a$,
$\mu = \lambda - 1$ ; l'invariance relie $h$ aux $a'_i$ par $c_i \langle x,
a'_i \rangle = \mu_i h(x, a_i)$, $c_i = h(a_i, a_i)$ réel. En posant $d_i =
h(a_{i-1}, a_i) = c_i/\mu_i$, la symétrie hermitienne donne la récurrence
\[
d_{i+1} = \bar\beta_i\, \bar d_i ,
\]
d'où $d_i$ comme produit alterné des $\beta_j$ et de leurs conjugués\footnote{La
ligne $d_3$ de la page 27 porte $d_0$ là où la récurrence donne $\bar d_0$ ;
la ligne générale $d_{2i+1}$ a bien $\bar d_0$. Les conventions de linéarité de
$h$ ne sont pas fixées sur la page ; les formules sont données ici à cette
réserve près.}. La normalisation $c_0 = 1$ et la réalité des $c_i$ imposent
des conditions de rationalité : si les $\beta_i$ sont dans $k_0$, les
$\mu_i/\mu_0$ ($i$ pair) et $\mu_i/\bar\mu_0$ ($i$ impair) doivent y être.
Page 29, en écrivant $\beta_i = \xi_i\mu_i\bar\mu_{i+1}$, il obtient $c_{i+1}
= N(\mu_{i+1})\,\xi_i\, c_i$ avec $\xi_i$ réel, d'où $c_n = N(\mu_1 \cdots \mu_n)
\xi_0 \cdots \xi_{n-1}$\footnote{Les indices de $N(\cdot)$ dans les deux
formules de droite de la page 29 semblent décalés d'un cran par rapport à la
liste de gauche ; la moitié supérieure de la page est barrée.}.

La page 28 calcule en rang $2$ : $\tau_0\tau_1$ a pour trace $\lambda_0 +
\lambda_1 + \beta$ et pour déterminant $\lambda_0\lambda_1$, et l'invariant de
conjugaison
\[
A(u) = \frac{(\operatorname{Tr} u)^2}{\det u} - 2
\]
vaut, pour $u$ de valeurs propres $e^{i\theta_1}, e^{i\theta_2}$, $2\cos(\theta_1 -
\theta_2)$. Il en dresse la liste des valeurs $-2, -1, 0, 1, \alpha, \alpha'$,
c'est-à-dire les $2\cos(2\pi k/p)$ pour $p = 2, 3, 4, 6$ et, avec $\alpha,
\alpha' = (-1 \pm \sqrt 5)/2$, pour $p = 5$ : les ordres de rotation dont la
trace est rationnelle ou dans $\mathbb{Q}(\sqrt 5)$\footnote{L'identification
des valeurs est de cette lecture. Un petit tableau de chiffres en marge ($1\,2\,2
/ 3\,3\,3 / 4\,4\,4 / 5\,5\,5$) et la ligne « $\tau_0\ \tau_1\ \tau_\infty$ /
$p\ q\ r$ », avec « un des trois $= 2$ déjà traité ; $p, q, r \geqslant 3$ »,
suggèrent une liste de triplets d'ordres, à la manière des groupes de
réflexions unitaires de rang $2$ ; la page n'en dit pas plus.}. La condition
$A(\tau_0\tau_1) = k$ donne une équation du second degré en $c/(2 - k_1)$, dont
il cherche les solutions réelles, et la positivité de la forme s'écrit $0 < c <
N(\mu_1) = 2 - k_1$, $k_1 = 2\cos(2\pi/\nu_1)$.

\pagerange{32}{35}
\section*{VIII. Détermination des $\tau_i$ par les rotations (pages 31 à 35)}

La garde porte « détermination de $(\tau_i)$ par $(\rho_i)$, i.e.\ d'une
représentation de $G_n$ via celle de $G_n^+$ ». Ici $\rho_i = \tau_i\tau_{i+1}$
désigne une rotation, non le continuant $\rho$\footnote{On garde la lettre de
la page dans cette seule section ; ailleurs $\rho$ est toujours $L_{n+1}$.}.

\emph{La donnée pseudo-sphérique} (page 32) : $(V, f_V, (L_i), a_0)$, $V$ de rang
$n+1$, $f_V$ quadratique, $L_i$ des droites et $a_0$ une base de $L_0$, avec
a) $f_V|_{L_i}$ lisse ; b) $L_i \perp L_j$ pour $i, j$ non consécutifs ; c)
$L_i \not\perp L_{i+1}$ sur chaque fibre ; d) $f_V(a_0) = 1$. Cela correspond aux
polytopes réguliers dont $\beta$ est inversible, sans que $f_V$ soit
nécessairement lisse. La normalisation $\varphi(a_i, a_{i+1}) = -f(a_{i+1})$
définit de proche en proche les bases $a_i$, et $a'_i = -\psi(a_i)/f(a_i)$ ; on
retrouve $f(a_i) = \beta_0 \cdots \beta_{i-1}$, $\varphi(a_i, a_{i+1}) = -\beta_0
\cdots \beta_i$, et
\[
\frac{\varphi(a_i, a_{i+1})^2}{f(a_i)\, f(a_{i+1})} = \beta_i = 4\cos^2\theta_i,
\qquad \alpha_i = \beta_i - 2 = 2\cos 2\theta_i ,
\]
où $\theta_i$ est l'angle des droites de $a_i$ et $a_{i+1}$ ; $\alpha_i$ est la
trace de la rotation $\tau_i\tau_{i+1}$ sur son plan.

\emph{La proposition marquée « ?? »} (page 33). Soient $G_2 = \langle \tau_0,
\tau_1, \tau_2 \mid \tau_i^2 = (\tau_0\tau_2)^2 = 1 \rangle$ et $G_2^+ =
\langle \rho_0, \rho_1 \mid (\rho_0\rho_1)^2 = 1 \rangle$ son sous-groupe des
rotations, $\rho_0 = \tau_0\tau_1$, $\rho_1 = \tau_1\tau_2$ (on a bien $\tau_0
= \rho_0\tau_1$, $\tau_2 = \tau_1\rho_1$, $\tau_0\tau_2 = \rho_0\rho_1$). Pour
$X$ fibré projectif de dimension relative $N \geqslant 1$, soit
$\operatorname{Hom}^!(G_2, \operatorname{Aut} X)$ l'ensemble des représentations
où les $\tau_i$ sont des réflexions et où $\rho_0^2$, $\rho_1^2$ et le
commutateur de $\rho_0, \rho_1$ sont non triviaux sur chaque fibre, et
$\operatorname{Hom}^!(G_2^+, \operatorname{Aut} X)$ celui des représentations
des rotations satisfaisant les mêmes conditions et, si $N \geqslant 2$,
l'existence locale de sous-fibrés $X_0, X_1, X_2$ de codimension $2$, d'intersection
de codimension $3$, sur lesquels $\rho_0, \rho_1, \rho_0\rho_1$ induisent
l'identité. \emph{Alors la restriction $\varphi \mapsto \varphi|_{G_2^+}$ serait
bijective.} Il marque l'énoncé « ?? » et « (?) », et ajoute page 34 : si $N = 2$,
il y a une unique conique invariante par $\varphi^+$, et elle l'est par
$\varphi$ --- « peut-être faut-il se donner en plus l'hypothèse conique lisse
invariante ». L'énoncé reste conjectural sur ces pages.

\emph{Ce qui est démontré} (page 35), dans le cadre orthogonal : si les $\beta_i$,
$2(\alpha_{i-1} + \alpha_i)$, $\alpha_{i-1} - 2$ et $\alpha_i - 2$ sont
inversibles, alors $\tau_i$ se reconstruit à partir de $\rho_{i-1} =
\tau_{i-1}\tau_i$ et $\rho_i = \tau_i\tau_{i+1}$ ($1 \leqslant i \leqslant
n-1$), et donc aussi $\tau_{i-1} = \rho_{i-1}\tau_i$, $\tau_{i+1} =
\tau_i\rho_i$. En effet $\rho_{i-1}$ fixe point par point $\{a_{i-1},
a_i\}^\perp$ et $\rho_i$ fixe $\{a_i, a_{i+1}\}^\perp$ ; les hypothèses disent
que $f$ est lisse sur les plans $\langle a_{i-1}, a_i \rangle$ (discriminant
proportionnel à $\beta_{i-1}(2 - \alpha_{i-1})$), $\langle a_i, a_{i+1} \rangle$
et sur l'espace $\langle a_{i-1}, a_i, a_{i+1} \rangle$ (demi-discriminant
proportionnel à $\alpha_{i-1} + \alpha_i$), et alors
\[
\bigl(V^{\rho_{i-1}} + V^{\rho_i}\bigr)^\perp = \langle a_{i-1}, a_i \rangle
\cap \langle a_i, a_{i+1} \rangle = \mathcal{O}_S\, a_i ,
\]
de sorte que $\tau_i$ est la réflexion orthogonale de direction cette
droite\footnote{En tête de page, les mêmes conditions de lissité ; ensuite la
page échange les noms : elle appelle $\alpha_i$ le quotient
$\varphi(a_i, a_{i+1})^2/f(a_i)f(a_{i+1}) = 4\cos^2\theta_i$, que la page 33
appelait $\beta_i$, et écrit $\beta_i = \alpha_i + 2 = 2(2\cos^2\theta_i + 1) =
4\cos 2\theta_i$, qui est faux : $2\cos^2\theta_i + 1 = \cos 2\theta_i + 2$. On
garde les noms de la page 33.}.

\pagerange{37}{42}
\section*{IX. $\mathcal{U}_\rho \simeq \operatorname{Aut}(X, (\tau_i))$ (pages 36 à 42)}

\emph{Le schéma en groupes} (page 37). Pour $\rho \in \Gamma(S, \mathcal{O}_S)$,
soit
\[
\mathcal{U}_\rho = \underline{\operatorname{Spec}}\, \mathcal{O}_S[T]_{(1 +
\rho T)}, \qquad \lambda *_\rho \lambda' = \lambda + \lambda' + \rho\lambda\lambda'
.
\]
C'est un schéma en groupes commutatif lisse de dimension relative $1$, et
$\lambda \mapsto 1 + \rho\lambda$ est un homomorphisme $\mathcal{U}_\rho \to
\mathbb{G}_m$, isomorphisme si et seulement si $\rho$ est inversible ; pour
$\rho = 0$, $\mathcal{U}_0 = \mathbb{G}_a$. C'est une famille qui déforme
$\mathbb{G}_m$ en $\mathbb{G}_a$\footnote{Ces schémas en groupes sont aujourd'hui
ceux de T.~Sekiguchi, F.~Oort et N.~Suwa (fin des années 1980), qui les
notent $\mathcal{G}^{(\lambda)}$ et en font le pont entre théorie de Kummer et
théorie d'Artin-Schreier ; les feuillets sont antérieurs. Cité de mémoire.}.

Pour un polytope régulier pseudo-sphérique d'invariants $\mu_i$ et $\beta_i$,
$\rho = \pi(\tilde c)$ est le continuant de la convention\footnote{La page
écrit $\rho = L_{n+1}(\mu_0, \lambda_0, \ldots, \lambda_{n-1}, \mu_n)$, les
arguments médians pouvant être des $\beta$ ; dans le cas des réflexions
c'est $L_{n+1}(\beta_0, \ldots, \beta_{n-1})$, comme on le montre aux pages 47
et 51 (section X).}, et
$\mathcal{U}_\rho$ opère fidèlement sur $E$ par $u_\lambda = \mathrm{id} +
\lambda\, \pi \otimes \tilde c$ (et sur le dual par la transposée), en
fixant $\pi$, $\tilde c$ et en commutant aux $\tau_i$ : la composition donne
bien $u_\lambda u_{\lambda'} = u_{\lambda *_\rho \lambda'}$ puisque $\pi(\tilde
c) = \rho$, et $\det u_\lambda = 1 + \rho\lambda$. D'où
\[
\mathbb{G}_m \times \mathcal{U}_\rho \longrightarrow
\underline{\operatorname{Aut}}(E, \varphi^E), \qquad \mathcal{U}_\rho
\longrightarrow \underline{\operatorname{Aut}}(X, \varphi^X),
\]
et le \emph{Théorème} de la page 38 : ces flèches sont des isomorphismes. De
plus, si $D_0 = H_1 \cap \cdots \cap H_n$ est la droite projective des points de
départ possibles, $d$ son point à l'infini ($D_0 \cap C$) et $D_0^{**} = D_0
\smallsetminus \{c, d\}$, alors $\mathcal{U}_\rho$ opère sur $D_0$ et $D_0^{**}$
est un torseur sous $\mathcal{U}_\rho$, trivialisé par $s_0$\footnote{Le
théorème est énoncé sans preuve. Le torseur se voit directement : les points
de $D_0$ sont les $s_0 + \lambda\tilde c$ et $\tilde c$, $d$ est le point où
$\pi(s_0 + \lambda\tilde c) = 1 + \rho\lambda$ s'annule, et $D_0^{**}$ est
exactement $\{s_0 + \lambda\tilde c : 1 + \rho\lambda \text{ inversible}\} =
\mathcal{U}_\rho \cdot s_0$. L'isomorphisme est repris, sous forme projective et
avec la même absence de preuve, aux pages 87 et 88.}.

\emph{La forme} (pages 39 à 42). Supposons les $\lambda_i = -1$ : il y a sur $E$
une forme invariante canonique $f_E$, unique par $f_E(s_0) = 0$, $f_E(a_0) =
1$, d'où une quadrique $Q_X \subset X$. Le calcul de la section VI donne que
$u_\lambda$ conserve $f_E$ si et seulement si $\lambda(\sigma +
\lambda\tilde\sigma\tilde\rho) = 0$. D'où la \emph{Proposition} de la page 42 :
\begin{enumerate}
  \item[1)] si $\sigma\tilde\rho$ est inversible, il existe une unique section
    de $\mathcal{U}_\rho$, partout $\neq 1$, qui respecte $f_E$ --- ou, ce qui
    revient au même, $Q_X$ : l'antipodisme ;
  \item[2)] sont équivalentes : a) $\sigma = 0$ et $\tilde\sigma\tilde\rho = 0$
    (la seconde condition suit de la première si $2$ est inversible) ; b)
    $\mathcal{U}_\rho$ respecte $f_E$ ; c) $\mathcal{U}_\rho$ respecte $Q_X$ ;
    et, si $\tilde\sigma\tilde\rho = 0$, d) il existe localement une section de
    $\mathcal{U}_\rho$ partout $\neq 1$ qui respecte $f_E$.
\end{enumerate}
Dans 1), « $\sigma\tilde\rho$ inversible » est exactement ce qu'il faut : pour
$n$ impair, $\sigma = 2\tilde\sigma$ contient l'inversibilité de $2$, sans
laquelle l'antipodisme est trivial (section VI)\footnote{Les pages 40 à 42
traitent séparément $n$ pair ($\lambda = -1/\tilde\rho$), $n$ impair
($\tilde\sigma(2 + \lambda\rho) = 0$, $u = \mathrm{id} - (2/\rho)\pi \otimes
\tilde c$) et l'exemple $n = 1$, où $\rho = 4 - \beta = 2 - \alpha$, $\sigma =
2$, $\tilde\sigma = 1$ : si $\rho$ est inversible, $\mathcal{U}_\rho =
\mathbb{G}_m$ et l'antipodisme est l'unique $u \neq 1$ respectant la conique
$Q_X$ ; $\mathcal{U}_\rho$ entier la respecte si et seulement si $\rho = 0$ en
caractéristique $2$. La prose de la page 41 est en bonne partie illisible.}.

\pagerange{44}{54}
\section*{X. Calculs de discriminants (pages 43 à 54)}

\subsection*{La matrice de la forme (pages 44 et 45)}

On suppose les $\lambda_i = -1$. La forme $f_E$ est caractérisée par : a) elle
est invariante par les $\tau_i$ ; b) $f_E(s_0) = 0$, donc $f_E$ s'annule sur tous
les sommets $s_i$ ; c) $f_E(a_0) = 1$. La quadrique $f_E = 0$ passe par les
sommets : c'est la \emph{sphère circonscrite}, prolongée à l'espace projectif.
Dans la base $(a_{-1}, a_0, \ldots, a_n)$,
\[
f_E(a_i) = c_i = \beta_0 \cdots \beta_{i-1}, \quad \varphi(a_i, a_{i+1}) =
-c_{i+1}, \quad \varphi(s_0, a_0) = -1, \quad \varphi(s_0, a_i) = 0 \ (i
\geqslant 1),
\]
les autres $\varphi(a_i, a_j)$ nuls, et $\varphi(s_i, s_j) = -c_i$ pour $i <
j$. La réflexion $\tau_i$ est orthogonale stricte si et seulement si $c_i$ est
inversible, et toutes le sont si et seulement si les $\beta_i$ le
sont\footnote{La page indexe $f(a_i) = c_{i-1}$, avec en marge « attention au
décalage des notations » ; on garde ici $c_i = f(a_i)$, comme les pages 15 à
17. Elle écrit « $\tau_i$ str.\ orth.\ $\Leftrightarrow (\alpha_0 + 2), \ldots,
(\alpha_{i-1} + 2) \in k^*$ », ce qui est la même chose.}.

\subsection*{Les déterminants tridiagonaux (pages 45, 46, 48 et 49)}

Les matrices de $\varphi$ sont tridiagonales. Notons $\Delta_m(C_0, \ldots,
C_{m-1})$ le déterminant de la matrice d'ordre $m$ de diagonale $(2C_0, \ldots,
2C_{m-1})$ et de coefficients voisins $-C_1, \ldots, -C_{m-1}$ ; il obéit à la
récurrence (encadrée page 48, déjà page 45)
\[
\Delta_m(C_0, \ldots, C_{m-1}) = 2C_0\, \Delta_{m-1}(C_1, \ldots, C_{m-1}) -
C_1^2\, \Delta_{m-2}(C_2, \ldots, C_{m-1}),
\]
avec $\Delta_1 = 2c_0$, $\Delta_2 = c_1(4c_0 - c_1)$, $\Delta_3 = 2c_2(4c_0c_1 -
c_0c_2 - c_1^2)$, valeurs justes\footnote{La page 48 écrit
$\Delta_{n+2}(0, 1, c_0, \ldots)$, déterminant de la forme sur $E$, égal au
déterminant d'ordre $n$ ; avec la tache qui peut couvrir un signe, c'est
$-\Delta_n$, comme à la page 45 où $\delta(f) = -F_n$.}. Pour un déterminant
tridiagonal quelconque (page 46), seules contribuent les permutations $\sigma$
telles que $\sigma(i) \in \{i - 1, i, i + 1\}$, c'est-à-dire les produits de
transpositions $(i, i+1)$ deux à deux disjointes. En substituant $c_i = \beta_0
\cdots \beta_i$, tout se factorise :
\[
\delta_n(\beta_0, \ldots, \beta_{n-1}) = \Delta_n(\beta_0, \beta_0\beta_1,
\ldots, \beta_0 \cdots \beta_{n-1}) = \beta_0^n \beta_1^{n-1} \cdots
\beta_{n-1}\, L_n(\beta_1, \ldots, \beta_{n-1}),
\]
où $L_k$ est le continuant des conventions, qui se développe sans
simplification (« il n'y a pas de cancellation de termes ») :
\[
L_k(B_1, \ldots, B_{k-1}) = \sum_{p \geqslant 0} (-1)^p\, 2^{k - 2p}
\sum_{J} B_J ,
\]
la seconde somme portant sur les parties $J \subset \{1, \ldots, k-1\}$ de
cardinal $p$ sans deux entiers consécutifs\footnote{Ce sont, à la normalisation
près, les \emph{continuants} d'Euler (le mot est de Sylvester), et le
développement est leur règle classique « des paires non adjacentes ». La page
écrit des bornes flottantes ($n$ au lieu de $n-1$ dans la dernière somme).
Page 49, $L_4$ est écrit avec un terme « $B_3B_4$ » qui n'existe pas pour trois
arguments ; le calcul qui suit garde exactement $B_1B_3$ et donne le résultat
juste, $L_4 = 16 - 4(B_1 + B_2 + B_3) + B_1B_3 = (\alpha_1 - 2)(\alpha_3 - 2) -
4(\alpha_2 + 2)$. Les deux premières lignes de la page 49 ($B_i = B'_i + 4$, un
produit de facteurs $2 - B_i/2$) sont des essais sans suite.} :
\[
L_2 = 4 - B_1, \qquad L_3 = 8 - 2(B_1 + B_2), \qquad L_4 = 16 - 4(B_1 + B_2 +
B_3) + B_1B_3 .
\]
Il note la symétrie $L_k(B_1, \ldots, B_{k-1}) = L_k(B_{k-1}, \ldots, B_1)$ et,
en bordant la matrice d'une ligne, $\delta_{n+1}(1, B_0, \ldots, B_{n-1}) =
B_0^n \cdots B_{n-1}\, L_{n+1}(B_0, \ldots, B_{n-1})$, avec les exemples justes
$\delta_2(1, B_0) = 4 - \alpha_0^2$, $\delta'_3(1, B_0, B_1) = -B_0^2B_1(\alpha_0 +
\alpha_1)$. La récurrence se transporte aux $L$ (page 51) :
\[
L_n(B_1, \ldots, B_{n-1}) = 2L_{n-1}(B_2, \ldots, B_{n-1}) - B_1\,
L_{n-2}(B_3, \ldots, B_{n-1}),
\]
et aux moitiés : $L'_{2m+1}(B_1, \ldots) = L_{2m}(B_2, \ldots) - B_1
L'_{2m-1}(B_3, \ldots)$, $L_{2m} = 4L'_{2m-1}(B_2, \ldots) - B_1 L_{2m-2}(B_3,
\ldots)$.

\emph{Modulo $2$} (page 50). La récurrence donne $L_k \equiv B_1 L_{k-2}(B_3,
\ldots) \pmod 2$, donc
\[
L_{2m}(B_1, \ldots, B_{2m-1}) \equiv B_1 B_3 \cdots B_{2m-1},
\]
\[
L'_{2m+1}(B_1, \ldots, B_{2m}) \equiv \sum_{j=0}^{m} (B_1 B_3 \cdots
B_{2j-1})(B_{2j+2} B_{2j+4} \cdots B_{2m}),
\]
chaque terme étant un produit d'indices impairs suivi d'un produit d'indices
pairs ; les termes $j = 0$ et $j = m$ sont $B_2B_4 \cdots B_{2m}$ et $B_1 B_3
\cdots B_{2m-1}$. Il en déduit les discriminants modulo $2$ ; celui de rang pair
est un carré.

\subsection*{Le centre, $\rho$ et $\sigma$ (pages 47 et 51)}

Le centre est le point fixe commun des $\tau_i$. On cherche $\tilde c =
\sum_{i=-1}^{n} \xi_i a_i$ avec $\langle \tilde c, a'_i \rangle = 0$ pour $0
\leqslant i \leqslant n$ et $\xi_n = 1$ ; les accouplements des conventions
donnent
\[
\xi_{n-1} = 2, \qquad \xi_{i-1} = 2\xi_i - \beta_i\, \xi_{i+1} ,
\]
d'où, par la récurrence des continuants et leur symétrie,
\[
\xi_{n-i} = L_i(\beta_{n-1}, \ldots, \beta_{n-i+1}) \quad (1 \leqslant i
\leqslant n+1),
\]
\[
\rho = \xi_{-1} = \pi(\tilde c) = L_{n+1}(\beta_0, \ldots, \beta_{n-1}), \qquad
\sigma = \xi_0 = L_n(\beta_1, \ldots, \beta_{n-1}).
\]
Les premières valeurs, calculées page 47 en $\alpha$, sont justes : $\xi_{n-2}
= 2 - \alpha_{n-1}$, $\xi_{n-3} = -2(\alpha_{n-1} + \alpha_{n-2})$, $\xi_{n-4} = -4 -
2(\alpha_{n-1} + 2\alpha_{n-2} + \alpha_{n-3}) + \alpha_{n-1}\alpha_{n-3}$.

On a $\varphi(a_i, \tilde c) = -c_i\langle \tilde c, a'_i \rangle = 0$ pour $i
\geqslant 0$ et $\varphi(s_0, \tilde c) = -\xi_0$, donc
\[
\psi(\tilde c) = -\sigma\pi, \qquad f_E(\tilde c) = \tfrac12 \langle \tilde c,
\psi(\tilde c) \rangle = -\tfrac12\sigma\rho = -\tilde\sigma\tilde\rho ,
\]
la dernière égalité ayant un sens sans diviser par $2$, puisque $\sigma\rho =
2\tilde\sigma\tilde\rho$ est une identité dans $\mathbb{Z}[\beta]$. Le centre est
à distance finie si et seulement si $\rho$ est inversible ; alors $c = \tilde
c/\rho$ et $f_E(c) = -\tilde\sigma\tilde\rho/\rho^2$.

On reconnaît dans $\rho$, à un facteur $2^{n+1}$ près, le déterminant de la
matrice de Gram $(-\cos(\pi/p_{ij}))$ du groupe de Coxeter : $\rho > 0$ pour un
polytope convexe (groupe fini), $\rho = 0$ pour un pavage euclidien, $\rho < 0$
dans le cas hyperbolique. Pour le carré, $\rho = 4 - 2 = 2$ ; pour le pavage
$\{4, 4\}$ du plan, $n = 2$, $\rho = 8 - 2(2 + 2) = 0$ ; pour $\{3, 6\}$, $\rho =
8 - 2(1 + 3) = 0$\footnote{C'est le critère de Schläfli ; le rapprochement est
de cette lecture. Le centre à l'infini quand $\rho = 0$ est le centre d'un
pavage euclidien, qui n'en a pas.}.

\subsection*{Discriminants et lissité (pages 51 et 52)}

La matrice de $\varphi$ dans la base $(a_i)$ de $V$ est le produit de la matrice
$(\langle a_j, a'_i \rangle)$, de déterminant $(-1)^{n+1}\rho$, par
$\operatorname{diag}(-c_i)$ ; d'où
\[
\delta(f_V) = \det \varphi_V = B\rho, \qquad \delta(f_E) = -B\sigma, \qquad B =
\beta_0^n \beta_1^{n-1} \cdots \beta_{n-1},
\]
et, pour les rangs impairs, les demi-discriminants $B\tilde\rho$, $B\tilde\sigma$.
Donc :
\begin{itemize}
  \item $f_E$ est lisse si et seulement si les $\beta_i$ et $\tilde\sigma$ sont
    inversibles --- si $2$ est inversible ou $n$ pair, si et seulement si les
    $\beta_i$ sont inversibles et $\psi(\tilde c)$ partout non nul ;
  \item $f_V$ est lisse si et seulement si les $\beta_i$ et $\tilde\rho$ sont
    inversibles --- si $2$ est inversible ou $n$ impair, si et seulement si les
    $\beta_i$ sont inversibles et $\pi(\tilde c) \neq 0$ sur chaque fibre :
    « condition à distance finie » ;
  \item $(E, V, f)$ est « birégulier » si et seulement si les $\beta_i$ et
    $f_E(\tilde c) = -\tilde\sigma\tilde\rho$ sont inversibles\footnote{Les
    discriminants de la page 52 sont notés $\delta_\varepsilon = B\sigma$,
    $\delta_V = B\rho$ (et $\delta'$ pour les moitiés), au signe près. Le mot
    « birégulier » est lu avec doute.}.
\end{itemize}

\subsection*{L'antipodisme et le plan $\langle s_0, \tilde c \rangle$ (pages 53 et 54)}

L'antipodisme est $\underline a = -\mathrm{id} + \pi \otimes c'$, défini si et
seulement si $\tilde\rho$ est inversible (c'est le cas si $f_V$ est lisse) ; on a
$\langle c', \pi \rangle = 2$, donc $-\underline a$ est une réflexion, sauf aux
points de caractéristique $2$ quand $n$ est impair, où $\underline a$ est
l'identité. Il commute aux $\tau_i$ et fixe $c'$. Dans $E$, ses points fixes
forment la droite de $c'$ si $2$ est inversible, et $V$ en caractéristique $2$ ;
dans $X$, $\mathbb{P}(V) \cup \{c\}$, réduit à $\mathbb{P}(V)$ en
caractéristique $2$. Il envoie $s_0$ sur $\tilde s_0 = c' - s_0$, le sommet
antipode.

Sur le plan engendré par $s_0$ et $\tilde c$, qui est la droite projective
$D_0$ des points de départ,
\[
g(\lambda, \mu) = f_E(\lambda s_0 + \mu\tilde c) = -\sigma\lambda\mu -
\tilde\sigma\tilde\rho\,\mu^2 .
\]
\begin{itemize}
  \item \emph{$n$ pair} : $g = -\mu\sigma(\lambda + \tilde\rho\mu)$. Elle est
    nulle si et seulement si $\sigma = 0$, et non nulle sur chaque fibre si et
    seulement si $\sigma$ est inversible ; cela équivaut à : $Q_X \cap D_0$ est
    un revêtement étale de degré $2$ de $S$, de sections $s_0$ et (quand
    l'antipodisme est défini) $\underline a\, s_0$ ; à $D_0 \not\subset Q_X$
    sur chaque fibre ; à $\psi(\tilde c) \neq 0$ sur chaque fibre. Les zéros
    avec $\mu$ inversible sont les multiples de $-\tilde\rho s_0 + \tilde c$,
    c'est-à-dire, si $\tilde\rho$ est inversible, de $c' - s_0 = \tilde s_0$.
  \item \emph{$2$ inversible}, $n$ quelconque : $g = -\sigma\mu(\lambda +
    \tfrac{\rho}{2}\mu)$, et tout ce qui précède vaut en remplaçant
    $\tilde\rho$ par $\rho/2$ ; le second zéro est $-s_0 + 2\tilde c/\rho =
    \underline a\, s_0$.
  \item \emph{$n$ impair, caractéristique $2$} : $g = \tilde\sigma\tilde\rho\,
    \mu^2$, identiquement nulle si et seulement si $\tilde\sigma\tilde\rho =
    0$, non nulle sur chaque fibre si et seulement si $\tilde\sigma$ et
    $\tilde\rho$ sont inversibles ; alors $D_0 \not\subset Q_X$, et $D_0$ est
    \emph{tangente} à $Q_X$ en $s_0$\footnote{La page écrit ce cas « B) $2
    \cdot 1_S = 0$ » sans préciser la parité ; pour $n$ pair, la formule du
    cas pair vaut en toute caractéristique, et c'est pour $n$ impair que
    $\sigma = 2\tilde\sigma$ s'annule. Le dernier alinéa de la page 54, sur les
    systèmes avec $n = 2m$ et $\sigma$ inversible, est barré.}.
\end{itemize}

\pagerange{56}{60}
\section*{XI. Calculs antérieurs (pages 55 à 60)}

La garde rose de la page 55 porte « Calculs antérieurs » (second mot lu avec
doute) ; les feuillets qui suivent reprennent en effet, sous une forme plus
brute, des calculs des sections précédentes.

\emph{Page 56.} Le composé $V \xrightarrow{\psi} V^\vee$ envoie $a_i$ sur $-c_i
a'_i$, et $a'_i = a^*_{i-1} - 2a^*_i + \beta_i a^*_{i+1}$ ; le déterminant de la
matrice tridiagonale d'ordre $n+1$ de diagonale $-2$ vaut $(-1)^{n+1}\rho$, le
facteur $B$ venant des $c_i$\footnote{La page écrit ce déterminant
« $= L_{n+1}(\beta_0, \ldots, \beta_{n-1}) = \rho$ », en omettant le signe
$(-1)^{n+1}$, qui vaut $1$ pour $n$ impair.}.

\emph{Page 57}, barrée d'une grande croix : un premier essai sur les
sous-variétés invariantes par les réflexions, en termes des $h_i$ et $H_i$ ---
contenues dans aucun $H_i$, dans un seul, la remarque que $Z \subset H_i$
entraîne $Z \subset H_{i-1}$. C'est l'ébauche du théorème des pages 79 à 82.

\emph{Page 58}, un demi-feuillet au crayon : la table des accouplements
$\langle a_i, a'_j \rangle$ ($1$ pour $j = i+1$, $\mu_i$ pour $j = i$, $\beta_j$
pour $j = i-1$, $0$ sinon), dessinée en graphe à deux rangées, et les
trivialisations $V \simeq \mathcal{O}_S^{n+1} \simeq V'$.

\emph{Page 60 : les bases duales.} Soit $a'_{n+1} \in E^\vee$ la forme de noyau
$\langle a_{-1}, \ldots, a_{n-1} \rangle$ normalisée par $\langle a_n,
a'_{n+1} \rangle = 1$, de sorte que $(a'_0, \ldots, a'_{n+1})$ est une base de
$E^\vee$. La base duale $(a^*_i)$ de $(a_i)_{-1 \leqslant i \leqslant n}$
s'exprime dans celle-ci par
\[
a^*_i = a'_{i+1} + 2a'_{i+2} + \rho_{i+1,i+1}\, a'_{i+3} + \cdots +
\rho_{i+1,n-1}\, a'_{n+1} \qquad (-1 \leqslant i \leqslant n),
\]
en particulier
\[
\pi = a^*_{-1} = a'_0 + 2a'_1 + \rho_{0,0}\, a'_2 + \cdots + \sigma^\vee a'_n +
\rho\, a'_{n+1}, \qquad
a^*_0 = -\psi(s_0) = a'_1 + 2a'_2 + \cdots + \gamma\, a'_n + \sigma\, a'_{n+1},
\]
avec $\rho_{0,0} = 4 - \beta_0 = 2 - \alpha_0$\footnote{On vérifie la formule
encadrée en l'évaluant sur les $a_k$ : le coefficient obtenu est $L_{k+1} -
2L_k + \beta_{k-1}L_{k-1}$, nul par la récurrence des continuants. La page a
deux glissements d'indice ($\rho_{2,2}\, a'_3$ pour $a'_4$ dans la ligne de
$a^*_1$, $2a_2$ pour $2a_0$ dans celle de $a'^*_2$) et écrit $\rho_{ij} =
L_{j-i+1}(\beta_j, \ldots, \beta_i)$ là où la symétrie demande $L_{j-i+2}$. Sa
lettre $\sigma'$ y est notre $\sigma^\vee$.}. Dualement, $\tilde c = a'^*_{n+1}
= a_n + 2a_{n-1} + \cdots + \sigma a_0 + \rho a_{-1}$, et $\langle \tilde c, \pi
\rangle = \rho$. La page rassemble enfin : $\psi(a_i) = -(\beta_0 \cdots
\beta_{i-1})\, a'_i$, $f(a_i) = \beta_0 \cdots \beta_{i-1}$, $f(s_0) = 0$,
$\varphi(a_i, a_{i+1}) = -\beta_0 \cdots \beta_i$, $f(\tilde c) =
-\tilde\sigma\tilde\rho$.

\pagerange{62}{70}
\section*{XII. Le système dual (pages 61 à 70)}

Les pages paires 62 à 70 sont de sa main ; les pages impaires 61 à 69 sont des
pièces dactylographiées sans rapport.

\emph{La forme duale et l'identité des continuants} (page 62). Le système dual,
sur $E^\vee$, a ses propres pseudo-réflexions, numérotées à rebours, et sa forme
invariante $f'$, avec $\psi'(a'_j) = -(\beta_{n-1} \cdots \beta_j)\, a_j$,
$f'(a'_j) = \beta_{n-1} \cdots \beta_j$, $f'(a'_{n+1}) = 0$, $\psi'(\pi) =
-\sigma^\vee \tilde c$. Alors
\[
\psi'\psi = \beta\,\mathrm{id}_E + \gamma\, \pi \otimes \tilde c, \qquad
\boxed{\sigma\sigma^\vee = \beta + \gamma\rho}.
\]
La première se vérifie sur la base : $\psi'\psi(a_i) = c_i (\beta_{n-1} \cdots
\beta_i)\, a_i = \beta a_i$, et $\psi'\psi(\tilde c) = \sigma\sigma^\vee \tilde
c$ ; la seconde en découle. C'est l'identité de Desnanot-Jacobi pour la matrice
tridiagonale qui définit $\rho$ : le produit des déterminants « coin supérieur »
et « coin inférieur » moins le produit du déterminant entier et du déterminant
central vaut le produit des coefficients hors diagonale\footnote{Le nom est de
cette lecture (on dit aussi « condensation de Dodgson »). La page 66 l'encadre
sous la forme $\sigma\sigma' - \beta = L_{n-1}(\beta_1, \ldots,
\beta_{n-2})\rho$, avec son $\sigma'$ pour notre $\sigma^\vee$ ; pour $n = 1$,
$4 - (4 - \beta_0) = \beta_0$.}.

\emph{L'inverse de $\psi$.} Posant $\Theta = \sigma\psi' + \gamma\,(\ell \mapsto
\ell(\tilde c)\,\tilde c)$, on a
\[
\Theta\psi = \beta\sigma\, \mathrm{id}_E ,
\]
puisque $\Theta\psi(x) = \sigma\beta x + \sigma\gamma\pi(x)\tilde c +
\gamma\varphi(x, \tilde c)\tilde c$ et $\varphi(x, \tilde c) = -\sigma\pi(x)$ :
$\beta\sigma$ est le « dénominateur » de $\psi^{-1}$\footnote{Les primes de
$\Theta, \Theta'$ sont croisées sur la page ; on nomme $\Theta$ l'opérateur
vérifié. Le bas de la page 62 est barré.}.

\emph{Les sommets du dual} (pages 62 à 66). Les vecteurs $s'_i$ du système dual
(sommets du polytope dual dans $E^\vee$) sont transportés dans $E$ par
$t_i = \Theta(s'_i)$ ; on a $t_i = \tau_i t_{i+1}$, $\langle t_i, \pi \rangle
= -\beta$, et
\[
t_{n+1} = \sigma\psi'(a'_{n+1}) + \gamma\,\tilde c, \qquad
t_i = t_{n+1} - \sigma\bigl[a_n + \beta_{n-1} a_{n-1} + \cdots + (\beta_{n-1}
\cdots \beta_i)\, a_i\bigr],
\]
$t_{n+1}$ n'ayant pas de composante sur $a_0$ (« $0 \cdot a_0$ ! ») et la
composante $-\beta$ sur $a_{-1}$. Comme $\psi(t_{n+1}) = \beta\sigma\,
a'_{n+1}$ et $\langle t_{n+1}, a'_{n+1} \rangle = \gamma$,
\[
f_E(t_{n+1}) = \tfrac12\beta\sigma\gamma = \beta\tilde\sigma\tilde\gamma ,
\qquad \varphi(t_{n+1}, a_n) = \beta\sigma.
\]
Si $\beta$ est inversible, $-t_i/\beta$ est un point de l'espace affine. Page 66,
avec $u = s_0 + s_1$ et $v = t_n + t_{n+1} = 2t_{n+1} - \sigma a_n$, il trouve
$f_E(u) = -1$ et $f_E(v) = \beta\sigma(2\gamma - \sigma)$, valeurs
justes\footnote{Sur la page, $\gamma$ et $\sigma'$ sont ceux des pages 60 à 70
($\gamma = \rho_{1,n-2}$, $\sigma' = \sigma^\vee$). Le haut des pages 62 et 66
est barré ; les coefficients de $t_0$ à la page 62 ne s'accordent pas avec
ceux de la ligne voisine, lus avec doute de part et d'autre.}.

\emph{Exemples} (page 68), qui vérifient tout ce qui précède. Pour $n = 1$ :
$\beta = \beta_0$, $\rho = 4 - \beta$, $\sigma = \sigma^\vee = 2$, $\gamma = 1$,
$t_2 = a_1 - \beta a_{-1}$, $f(t_2) = \beta$. Pour $n = 2$ : $\rho = -2(\alpha_0
+ \alpha_1)$, $\sigma = 2 - \alpha_1$, $\sigma^\vee = 2 - \alpha_0$, $\gamma = 2$,
$\Theta = (2 - \alpha_1)\psi' + 2\,\tilde c \otimes \tilde c$, $t_3 = 2a_2 +
\beta_1 a_1 - \beta_0\beta_1 a_{-1}$, $f(t_3) = (4 - \alpha_1^2)(2 + \alpha_0)$.

\emph{Le centre à distance finie} (page 70). Si $\rho$ est inversible, $E
\simeq V \oplus \mathcal{O}_S \tilde c$, et les projections $p_\lambda(x) = x/\lambda
- c$ ramènent les points de l'hyperplan $\pi = \lambda$ dans $V$, l'origine au
centre. Le sommet recentré est
\[
-\rho\,\tilde s_0 = \sigma a_0 + \rho_{n-1,2}\, a_1 + \cdots + \rho_{n-1,n-1}\,
a_{n-2} + 2a_{n-1} + a_n ,
\]
c'est-à-dire $-\rho\tilde s_0 = \sum_{i \geqslant 0}\xi_i a_i$ ; il calcule
de même $\tilde u = \tfrac12 u - c = \tilde s_0 + \tfrac12 a_0$ et les
$\tilde t_i$, et obtient
\[
-\frac{\rho\beta}{\sigma}\,\tilde t_{n+1} = \sigma^\vee \pi - \rho\, a^*_n
\]
dans ses notations\footnote{Cette page est la moins sûre de la section : ses
coefficients de $a_{n-1}$ et $a_n$ ont des signes surchargés qui ne suivent pas
tous le calcul (le calcul juste donne $-\sigma(2\sigma^\vee - \rho)/\beta\rho$
pour $a_{n-1}$), elle applique $p_{-\beta}$ en l'appelant $p_\beta$, et la
dernière ligne porte $+\rho\beta(\ldots)$ là où $t_i = t_{n+1} - \sigma[\ldots]$
donne $-\rho[\ldots]$. On ne reproduit que ce qui se vérifie.}.

\pagerange{72}{77}
\section*{XIII. Calculs centrés au centre, si celui-ci est à distance finie (pages 71 à 77)}

La garde porte « Calculs centrés en $\sigma$ (si centre à dist.\ finie, i.e.\
$\rho$ inv.) », la lettre $\sigma$ y désignant le centre.

\emph{Reformulation} (pages 72 et 73). Si $\rho$ est inversible --- $n$ impair
et $\tilde\rho$ inversible, ou $n$ pair, $\tilde\rho$ et $2$ inversibles ---, on
prend le centre pour origine et l'on identifie l'espace affine à $V$. La donnée
devient : A) $\varphi_G \colon \mathfrak{G}_n \to \operatorname{GL}(V)$ avec
$1^\circ)$ les $\tau_i^V$ pseudo-réflexions, $2^\circ)$ $V = \bigoplus L_i$, ou
encore $\bigcap H_i = 0$ sur chaque fibre, $3^\circ)$ $L_i \not\subset H_{i+1}$
sur chaque fibre\footnote{La page écrit « $L_i \subset H_{i+1}$ chaque fibre » ;
c'est la condition de non-dégénérescence 4) de la page 14, $L_i \not\subset
H_{i+1}$, qui est entendue.} ; B) le sommet recentré $\tilde s_0$, base de $D_0
= \bigcap_{i \geqslant 1} H_i$, qui fixe une base $a'_0$ de $L'_0$ par $\langle
\tilde s_0, a'_0 \rangle = 1$, puis $a_0$ par $\tau_0 = \mathrm{id} + a'_0
\otimes a_0$, puis les $a_i, a'_i$ de proche en proche. Pour des réflexions, le
sommet est
\[
\tilde s_0 = -\frac{1}{\rho}\,(\xi_0 a_0 + \xi_1 a_1 + \cdots + \xi_n a_n),
\qquad \xi_n = 1,\ \xi_{n-1} = 2,\ \xi_i = L_{n-i}(\beta_{i+1}, \ldots,
\beta_{n-1}),\ \xi_0 = \sigma. \quad (*)
\]
C) Les formes quadratiques invariantes forment un module localement libre de
rang $1$, normalisé par $f_V(a_0) = 1$.

\emph{Discriminant et sphère circonscrite} (page 74). Avec $\omega_V = a_0 \wedge
\cdots \wedge a_n$,
\[
\delta(f_V) = B\rho\; \omega_V^{\otimes -2},
\]
et, $\rho$ étant ici inversible, $f_V$ est lisse si et seulement si les
$\beta_i$ le sont. Posant
\[
\lambda = f_V(\tilde s_0) = \frac{\tilde\sigma\tilde\rho}{\rho^2} =
\begin{cases} \dfrac{\tilde\sigma}{4\tilde\rho} & n \text{ pair}, \\[1.5ex]
\dfrac{\tilde\sigma}{\tilde\rho} & n \text{ impair}, \end{cases}
\]
la sphère circonscrite au polytope est la quadrique $f_V = \lambda$ ; elle est
lisse si de plus $\tilde\sigma$ est inversible. Le calcul de la page 77
retrouve $\lambda$ : $f_V(\tilde s_0) = f(s_0) + f(c) - \varphi(s_0, c) =
-\tilde\sigma\tilde\rho/\rho^2 + \sigma/\rho$. Vérification sur deux cas : pour le
$p$-gone régulier d'arête $1$ ($n = 1$), $\lambda = 1/(4 - \beta_0) =
1/(4\sin^2(\pi/p))$, le carré du rayon circonscrit ; pour le cube d'arête $1$
($n = 2$, $\beta = (2, 1)$, $\tilde\sigma = 3$, $\tilde\rho = 1$), $\lambda =
3/4$\footnote{Ces deux vérifications sont de cette lecture.}.

\emph{Partir des réflexions orthogonales} (pages 74 à 76). (D) Soient donnés
$V$, $f_V$ et $\varphi^V \colon \mathfrak{G}_n \to \mathrm{O}(f_V)$ envoyant les
$\tau_i$ sur des réflexions orthogonales et satisfaisant $2^\circ)$, $3^\circ)$,
avec $f|_{L_0}$ partout non nulle. Cela provient d'une situation comme ci-dessus :
il suffit de choisir $a_0 \in L_0$ avec $f_V(a_0) = 1$, ce qui n'est possible
que localement pour une topologie où l'on peut extraire des racines carrées, le
choix étant fixé au signe près --- l'indétermination est un torseur sous
$\mu_2$, comme il le dit page 75\footnote{La page dit « il suffit de choisir une
section $a_0$ de $D_0$ telle que $f_V(a_0) = 1$ », puis précise
l'indétermination ; la remarque sur la topologie (étale si $2$ est inversible,
plate en général) est de cette lecture. « Rien n'empêche d'ailleurs que $f_V(s_0)
= 0$ », ajoute-t-il : le sommet peut être sur le cône isotrope.}. (E) Si l'on se
donne seulement $\varphi^C \colon \mathfrak{G}_n \to \operatorname{Aut}(C, Q_C)$
avec les conditions 1) à 3), cela provient localement d'un $(V, f_V,
(\tau_i^V))$ comme en (D), $V$ déterminé à un $\otimes L$ près, $f_V$ à une
unité près, et l'on rigidifie par $a_0$. D'où l'énoncé de la page 76 :
\emph{la catégorie des réalisations géométriques non dégénérées à centre à
distance finie est équivalente à celle des triples $(C, Q_C, \varphi^C)$ d'un
fibré projectif de dimension relative $n$, d'une hyperquadrique et d'une
représentation satisfaisant $1^\circ)$ à $3^\circ)$} --- énoncé sans
démonstration, et entendu aux rigidifications près qu'on vient de dire. Pour $n
= 2$, $Q_C$ est redondante : c'est l'unique conique lisse invariante\footnote{La
page écrit « Pour $n \neq 2$ » et met « NB $n \neq 2$ » en marge de l'énoncé ; la
marge de la page 75 explique que pour $n = 1$ une réflexion de $C =
\mathbb{P}(V)$ ne provient pas d'une réflexion de $V$, et la parenthèse finale de
la page 76 dit que pour $n = 2$ la conique se déduit canoniquement. Le
relèvement exige $n \geqslant 2$ (section IV) ; le cas $n = 2$ n'est pas
exclu, il est seulement plus simple.}.

\emph{La sphère médiane} (page 77). Avec $s'_i = -\rho\tilde s_i$, on a
$f_V(s'_0) = \tilde\rho\tilde\sigma$, $\varphi_V(s'_0, a_0) = \rho$,
$\varphi_V(s'_0, a_i) = 0$ pour $i \geqslant 1$, $s'_1 = s'_0 - \rho a_0$, et
\[
f_V(s'_0 + s'_1) = 4\tilde\rho\tilde\sigma - \rho^2 =
\begin{cases} 4\tilde\rho(\tilde\sigma - \tilde\rho) & n \text{ pair}, \\
\tilde\rho(4\tilde\sigma - \tilde\rho) & n \text{ impair}. \end{cases}
\]
Le point $(s'_0 + s'_1)/2$ est le milieu d'une arête. L'homothétie de rapport
$\lambda$ qui envoie la sphère circonscrite sur la sphère passant par les
$s'_0 + s'_1$ est donc donnée par
\[
\lambda = 2\sqrt{1 - \tilde\rho/\tilde\sigma} \quad (n \text{ pair}), \qquad
\lambda = 2\sqrt{1 - \tilde\rho/4\tilde\sigma} \quad (n \text{ impair}),
\]
et $\lambda/2$ est le rapport du rayon de la sphère médiane --- tangente aux
arêtes en leurs milieux --- au rayon circonscrit. Pour le cube, $(\lambda/2)^2
= 2/3$, ce qui est juste\footnote{L'interprétation par la sphère médiane et la
vérification sont de cette lecture ; la page parle de « la sphère circonscrite
aux transformés des $s'_0 + s'_1$ ».}.

\pagerange{79}{82}
\section*{XIV. Sous-espaces invariants (pages 78 à 82)}

On est sur un corps $k$, dans une réalisation projective non dégénérée : $c$
est l'intersection des $H_i$, $C$ l'hyperplan engendré par les $h_i$. Pour $-1
\leqslant i \leqslant n$, posons
\[
Z'_i = H_0 \cap \cdots \cap H_i \ (\dim n - i), \qquad Z^*_i =
\operatorname{Lin}(h_{i+1}, \ldots, h_n) \ (\dim n - i - 1),
\]
avec $Z'_{-1} = X$, $Z^*_{-1} = C$, $Z'_n = \{c\}$, $Z^*_n = \emptyset$. On a
toujours les deux chaînes $Z^*_n \subset \cdots \subset Z^*_{-1}$ et $Z'_n
\subset \cdots \subset Z'_{-1}$, et, pour $0 \leqslant i \leqslant n-1$, $Z^*_i
\subset Z'_i$, alors de codimension $1$, si et seulement si $h_{i+1} \in H_i$, c'est-à-dire $\beta_i = 0$ --- les autres
$h_j$ ($j \geqslant i + 2$) étant dans $H_0, \ldots, H_i$ par la
commutation\footnote{La page 81 écrit « $Z_i \subset Z'_i$, ce qui équivaut
aussi à $h_{i+1} \notin H_i$ i.e.\ $\beta_i = 0$ » ; c'est $\in$ qu'il faut,
$\beta_i = \langle a_{i+1}, a'_i \rangle$ mesurant précisément l'appartenance de
$h_{i+1}$ à $H_i$. La note de marge de la page 79 a la bonne équivalence. Les
deux définitions de la page 79 échangent $Z^*_{-1}$ et $Z'_{-1}$ par rapport à
la chaîne qui suit ; on suit la chaîne.}.

\emph{Théorème} (page 79). Un sous-espace projectif $Z \subset X$ est invariant
par $\mathfrak{G}_n$ si et seulement si ou bien $Z$ est l'un de $\emptyset$,
$X$, $\{c\}$, $C$ ; ou bien il existe $0 \leqslant i \leqslant n-1$ avec $\beta_i
= 0$, et $Z = Z'_i$ ou $Z = Z^*_i$.

\emph{Démonstration} (pages 81 et 82). $Z$ est invariant par une
pseudo-réflexion de centre $h$ et cocentre $H$ si et seulement si $h \in Z$ ou $Z
\subset H$. Si $Z \subset H_i$ avec $i > 0$, alors $Z \subset H_{i-1}$ : sinon
$h_{i-1} \in Z \subset H_i$, contre $h_{i-1} \notin H_i$. De même, si $h_j \in Z$
avec $j < n$, alors $h_{j+1} \in Z$ : sinon $Z \subset H_{j+1}$ et $h_j \in
H_{j+1}$. L'ensemble des $i$ tels que $Z \subset H_i$ est donc un segment
initial $[0, i_0]$, et pour $j > i_0$ on a $h_j \in Z$. Si $i_0 = n$, $Z \subset
\{c\}$ ; si aucun, $Z \supset C$ et $Z = C$ ou $X$ ; sinon $Z^*_{i_0} \subset Z
\subset Z'_{i_0}$, ce qui force $\beta_{i_0} = 0$ et, pour une raison de
dimension, $Z = Z^*_{i_0}$ ou $Z'_{i_0}$. La réciproque est claire. La marge de
la page 81 en tire, pour $i$ fixé, l'équivalence de : a) $Z^*_i \subset Z'_i$ ; b)
$Z^*_i$ invariant ; c) $Z'_i$ invariant ; d) $\beta_i = 0$.

\emph{Comptes} (pages 79 et 80). Pour chaque dimension $0 \leqslant d \leqslant
n$, il y a exactement deux espaces des chaînes, $Z'_{n-d}$ et $Z^*_{n-d-1}$, et
il note qu'ils coïncident exactement quand, pour $i = n - d$, $\mu_i = 0$ (ce qui,
pour des réflexions, est la caractéristique $2$) et $\beta_i = \beta_{i-1} = 0$,
la condition sur $\beta_i$ disparaissant pour $i = n$ et celle sur $\beta_{i-1}$
pour $i = 0$ ; si $\rho \neq 0$ ils sont tous distincts. D'où les corollaires :
\begin{enumerate}
  \item[(b)] il y a au plus $2(n+2)$ sous-espaces invariants, et exactement
    $2(n+2)$ si et seulement si tous les $\beta_i$ sont nuls et tous les $\mu_i$
    non nuls ;
  \item[(c)] il y a toujours $\emptyset$, $X$, $\{c\}$, $C$ ; pour que ce soient
    les seuls, il suffit que tous les $\beta_i$ soient non nuls\footnote{La page
    poursuit par « il faut et suffit » avec une condition plus faible, qui
    épargne $\beta_0$ quand $\mu_0 = 0$ et $\beta_{n-1}$ quand $\mu_n = 0$ (en
    caractéristique $2$, seuls les $\beta_i$ intérieurs seraient requis). Les
    deux énoncés se suivent sans que le premier soit barré. Tel que le théorème
    est écrit, $\beta_{n-1} = 0$ fournit l'espace invariant $Z'_{n-1}$, droite
    passant par $c$, qui n'est aucun des quatre dès que $n \geqslant 2$ ; la
    condition faible ne suit donc pas de l'énoncé, et cette lecture ne la
    tranche pas.}.
\end{enumerate}

\emph{Irréductibilité sur $C$} (pages 80 à 82). Les sous-espaces invariants
contenus dans $C$ sont $\emptyset$, $C$, les $Z^*_i$ pour $\beta_i = 0$, et
$\{c\}$ si $c \in C$, c'est-à-dire si $\rho = 0$. Donc \emph{$\mathfrak{G}_n$
opère irréductiblement sur $C$ si et seulement si tous les $\beta_i$ sont non nuls
et $\rho \neq 0$}, c'est-à-dire si et seulement si la forme bilinéaire $\varphi_V$
est non dégénérée, puisque $\det \varphi_V = B\rho$ ; pour $n$ pair cela exige la
caractéristique $\neq 2$, $\rho = 2\tilde\rho$ étant alors nul\footnote{Le d)
barré de la page 80 dit exactement cela (« il faut et suffit que $\rho \neq 0$ et
les $\beta_i \neq 0$ ») ; le c) de la page 82 omet $\rho \neq 0$, et attache la
condition de caractéristique à $n$ impair. Le d) repris page 80, qui décrit les
sous-espaces stables de $C$ comme les $Z^*_i$ pour $\beta_i = 0$, porte en marge
« si $c \notin C$ i.e.\ $\rho \neq 0$ », et s'arrête sur « Ces sous- ».}. La
\emph{Remarque} de la page 81 précise, pour $n \geqslant 2$ : il n'y a pas
d'hyperplan de $C$ invariant si et seulement si $\beta_0 \neq 0$ et l'on n'a pas
à la fois $\beta_1 = 0$ et $H_0 \cap H_1 \subset C$.

\pagerange{84}{94}
\section*{XV. Groupes engendrés par des pseudo-réflexions projectives, et réalisations géométriques (pages 83 à 94)}

\subsection*{Réalisations d'une carte régulière (pages 84 à 86)}

Soit $C$ une $n$-carte régulière combinatoire, $R$ l'ensemble de ses repères,
torseur à gauche sous $G = \operatorname{Aut} C$, le commutant de $G$ étant
formé des opérations du groupe cartographique, engendré par $\sigma_0, \ldots,
\sigma_n$\footnote{Ici $C$ désigne la carte, comme sur la page, et non
l'hyperplan à l'infini ; $\sigma_i$ sont les générateurs du groupe
cartographique, comme dans les dossiers 76 et 88, non le continuant.}. Un repère
$r$ définit des générateurs $\tau_i^{(r)}$ de $G$ par $\tau_i^{(r)} r = \sigma_i
r$. Une réalisation géométrique projective de $C$ dans un fibré projectif $X$
de dimension relative $n+1$ est la donnée d'une action $\varphi_X \colon G \to
\operatorname{Aut}_S X$ et d'un drapeau $X_0 \subset X_1 \subset \cdots \subset
X_n$ avec $\tau_i X_j = X_j$ pour $i \neq j$, les $\tau_i$ étant des
réflexions ; cela définit $R \to \mathrm{Drapeaux}(X)$.

\emph{Proposition} (pages 84 et 85, reprise page 87). Sur un corps, sont
équivalentes :
\begin{itemize}
  \item a) l'image des sommets engendre linéairement $X$ ; a') l'intersection
    des hyperplans images des facettes est vide ;
  \item b) aucune des applications $\Phi_i \to \mathrm{Grass}_i(X)$ (faces de
    dimension $i$ vers sous-espaces) n'est constante\footnote{La page écrit
    « $\mathrm{Gram}_i(X)$ » ; on attend une grassmannienne.} ;
  \item c) $\tau_i X_i \neq X_i$ pour tout $i$ ;
  \item d) deux faces distinctes de dimension $i$ entre une face de dimension
    $i-1$ et une de dimension $i+1$ ont des images distinctes ;
  \item e) les points $s_0$ (le point $X_0$), $s_{i+1} = \tau_i s_i$ sont
    projectivement indépendants ; e') la condition duale ;
  \item f) les centres $h_i$ et cocentres $H_i$ vérifient : $\alpha$) les $h_i$
    sont indépendants (ils engendrent un hyperplan $C$) ; $\alpha'$) les $H_i$
    sont indépendants ($\bigcap H_i = \{c\}$) ; $\beta$) $h_i \notin H_{i+1}$ ;
    et le point de départ vérifie $\delta$) $s_0 \in D_0 = H_1 \cap \cdots
    \cap H_n$, $s_0 \neq c$, $s_0 \neq d = D_0 \cap C$.
\end{itemize}
Les implications a) $\Rightarrow$ b), a') $\Rightarrow$ b), b) $\Leftrightarrow$
c) $\Leftrightarrow$ d), e) $\Rightarrow$ a), e) $\Rightarrow$ a') sont
triviales ; la seule qui ne le soit pas est c)
$\Rightarrow$ e), et la preuve en donne davantage :

\emph{Corollaire} (F). $X_i = \operatorname{Lin}(s_0, \ldots, s_i)$ pour $0
\leqslant i \leqslant n+1$ ($X_{n+1} = X$). Par récurrence : $s_i \in X_i
\subset X_{i+1}$ et $\tau_i X_{i+1} = X_{i+1}$ donnent $s_{i+1} \in X_{i+1}$ ;
si l'on avait $s_{i+1} \in X_i$, alors $\tau_i X_i = \operatorname{Lin}(\tau_i
X_{i-1}, \tau_i s_i) = \operatorname{Lin}(X_{i-1}, s_{i+1}) \subset X_i$,
contre c). Donc $s_{i+1} \notin X_i$, et $X_{i+1} = \operatorname{Lin}(X_i,
s_{i+1})$\footnote{La page écrit deux fois « $\in$ » (« Reste à prouver que
$s_{i+1} \in X_i$, i.e.\ $\tau_i s_i \in X_i$ ») là où le raisonnement demande
$\notin$.}. On note aussi (5) : $\tau_j s_i = s_i$ pour $j > i$, puisque $\tau_j$
commute à $\tau_{i-1}, \ldots, \tau_0$ et fixe $s_0$.

La page 86, sans phrases, dresse le treillis de la situation : $c =
\bigwedge H_i$, $C = \bigvee h_i$, $Y_{i-1} = \bigvee_{j < i} h_j$, $Z_{i+1} =
\bigwedge_{j > i} H_j$, $d = Z_1 \wedge C$, $D = Y_{n-1} \vee c$, et les
formules
\[
X_i = X_n \wedge Z_{i+1} = X_0 \vee Y_{i-1}, \qquad X'_i = Y_{i-1} \vee c =
Z_{i+1} \wedge C, \qquad X''_i = Y_{i-1} \vee d = Z_{i+1} \wedge D ,
\]
avec la dualité $h_i \leftrightarrow H_i$, $c \leftrightarrow C$, $d
\leftrightarrow D$.

\subsection*{Le torseur des drapeaux (pages 87 à 90)}

\emph{B)} Pour $(\tau_i)$ donné, il existe localement un drapeau satisfaisant
les conditions équivalentes si et seulement si les $\tau_i$ vérifient
$\alpha$), $\alpha'$), $\beta$). \emph{C)} Sous ces conditions, $(X_i)
\mapsto s_0$ est une bijection de l'ensemble des drapeaux non dégénérés
adaptés aux $\tau_i$ sur l'ensemble des sections de $D_0^{**} = D_0
\smallsetminus \{c, d\}$ ; plus généralement $(X_*) \mapsto X_{i_0}$ est une
bijection sur les sections de $D_{i_0}^{**}$, où $D_i$ est le fibré des
sous-espaces de dimension $i$ compris entre $Y_{i-1}$ et $Z_{i+1}$, privé des
deux sections $X'_i$, $X''_i$. En marge : les $X_i$ se déterminent les uns les
autres, $X_i = X_j \wedge H_j \wedge \cdots \wedge H_{i+1}$ et $X_j = X_i \vee
h_i \vee \cdots \vee h_{j-1}$ pour $i < j$.

\emph{D)} Soit $\underline{g}$ le schéma en groupes des automorphismes de $(X,
(\tau_i))$. Il opère sur $(D_0, c, d)$, et l'homomorphisme
\[
\underline{g} \longrightarrow \underline{\operatorname{Aut}}(D_0;\, c + d),
\]
vers les automorphismes de la droite $D_0$ qui induisent l'identité sur le
diviseur $c + d$, est un isomorphisme ; $D_0^{**}$ est un torseur sous
$\underline{g}$, et le schéma des drapeaux non dégénérés aussi. $\underline g$
est lisse, commutatif, de dimension relative $1$, de fibres $\mathbb{G}_m$ là
où $c \neq d$ et $\mathbb{G}_a$ là où $c = d$\footnote{« Induisent l'identité
sur $c + d$ » rend ses annotations « id sur », « produit de div.\ relatifs ».
Quand $c = d$, fixer seulement le point donnerait le groupe affine, de
dimension $2$ ; c'est l'identité sur le diviseur double qui donne
$\mathbb{G}_a$. Ce groupe est le $\mathcal{U}_\rho$ de la section IX, $c = d$
correspondant à $\rho = 0$ ; le rapprochement est de cette lecture.}.

\emph{Scholie} (page 88). Un système $(\tau_*, X_*)$ non dégénéré est une
\emph{rigidification} du système $(\tau_*)$. La catégorie des $(\tau_*)$ non
dégénérés équivaut à celle des $(\tau_*, X_*)$ munis d'un torseur sous
$\underline{g}$. Si le centre est partout à distance finie, en prenant $s_0$
comme $1$, $c$ comme $0$ et $d$ comme $\infty$ sur $D_0$, ce groupe est
$\mathbb{G}_m$, et la catégorie des $(\tau_*)$ non dégénérés équivaut au produit
de celle des $(\tau_*, X_*)$ par celle des modules inversibles sur $S$.

\emph{E)} (pages 88 à 90). Pour un quasi-polyèdre combinatoire régulier $\Pi$,
de groupe $G$ et de torseur des repères $R$, le choix de $r \in R$ définit un
épimorphisme $\mathfrak{G}_n \to G$ ; changer $r$ en $gr$ conjugue le système par
$\operatorname{int}(g)$. Les conditions (être formé de pseudo-réflexions, être
non dégénéré) ne dépendent pas de $r$, et, $G$ étant simplement transitif sur
$R$, tout objet attaché intrinsèquement à un système non dégénéré --- $D_0$, $c$,
$d$ --- l'est à $\varphi_G$. On trouve $\underline{\operatorname{Aut}}(X,
\varphi_G) \simeq \underline{\operatorname{Aut}}(D_0;\, c + d)$, et la catégorie
des $\varphi_G$ quasi-polyédrales non dégénérées équivaut à celle des couples
formés d'une réalisation géométrique projective non dégénérée de $\Pi$ et d'un
torseur sous $\underline{\operatorname{Aut}}(D_0;\, c + d)$. Si $c = d$,
$D_0 \smallsetminus \{c\}$ est un fibré affine de rang $1$ trivialisé par $s_0$,
et le groupe est le fibré vectoriel de ses translations.

\subsection*{Coordonnées universelles (pages 91 à 93)}

On prend $C$ pour hyperplan à l'infini ; $s_0, \ldots, s_{n+1}$ forment une base
affine de $X \smallsetminus C$, et l'enveloppe vectorielle s'identifie à
\[
\mathcal{E} = \mathcal{O}_S s_0 \oplus \cdots \oplus \mathcal{O}_S s_{n+1}
\xrightarrow{\ \pi\ } \mathcal{O}_S, \qquad \textstyle\sum \lambda_i s_i
\mapsto \sum \lambda_i .
\]
Les pseudo-réflexions $\tau_i = \mathrm{id} + a'_i \otimes a_i$ doivent
fixer les $s_k$ pour $k < i$ (la remarque (5) de la page 85), ce qui donne à
$a'_i$ la forme triangulaire $a'_i = \sum_{k \geqslant i} \lambda_{ik} s'_k$
dans la base duale $(s'_k)$ ; vérifier (2) $s_{i+1} = \tau_i s_i$ ; et (1) la
moitié restante des conditions de commutation, $\langle a_j, a'_i \rangle = 0$
pour $j \geqslant i + 2$\footnote{Les indices de la condition (1), entre
crochets sur la page, sont surchargés ; on donne la condition que le calcul
utilise.}. (2) donne $\lambda_{ii} a_i = s_{i+1} - s_i$, et l'on normalise
$\lambda_{ii} = 1$ :
\[
a_i = s_{i+1} - s_i, \qquad a'_i = s'_i + \lambda_i s'_{i+1} + \mu'_i(s'_{i+2} +
\cdots + s'_{n+1}) \quad (i < n), \qquad a'_n = s'_n + \lambda_n s'_{n+1} ,
\]
la condition (1) imposant l'égalité des coefficients au-delà de
$s'_{i+1}$\footnote{La lettre devant la dernière parenthèse est surchargée et lue
$\mu_i$ ; on la note $\mu'_i$ pour ne pas la confondre avec $\mu_i = \lambda_i
- 1$. Les deux formules de $\langle a_{i+1}, a'_i \rangle$ de la page 93
($1 + \alpha_i - \lambda_i$ dans la matrice, $\alpha_i - \lambda_i$ sous
l'accolade) ne s'accordent pas, et la lettre $\alpha_i$ n'y est pas définie ;
avec les formules ci-dessus, $\langle a_{i+1}, a'_i \rangle = \mu'_i - \lambda_i$
et $\beta_i = \mu'_i - \lambda_i$.}. Alors $\langle a_{i-1}, a'_i \rangle = 1$,
$\langle a_i, a'_i \rangle = \lambda_i - 1$ : $\lambda_i$ est la valeur propre de
$\tau_i$, et $\tau_i^2 = 1$ si et seulement si $\lambda_i = -1$ (page 93, juste).
Les $2n+1$ paramètres $(\lambda_i, \beta_i)$ sont les coordonnées universelles ;
avec le théorème de la section suivante, ils donnent la classification
annoncée à la page 13\footnote{Cette conséquence est de cette lecture : dans
ces coordonnées les conditions (B) du théorème sont satisfaites pour tous
paramètres ($\nu_n = 1$), et l'on obtient une réalisation non dégénérée pour
toute valeur des $\beta_i$ et toute valeur inversible des $\lambda_i$.}.

La page 92, barrée et dont la place dans la suite des feuillets n'est pas sûre,
renvoie à des formules (5), (8), (9) et à des sous-modules $F_n$ qui sont ceux du
théorème des pages 96 à 99 : c'en est un premier état. Elle esquisse un
« Corollaire » et un « Re Corollaire » pour un fibré vectoriel et des conditions
fibre à fibre. La page 93 calcule $\tau_i\tau_{i+1} = \mathrm{id} + a'_i \otimes
a_i + a'_{i+1} \otimes a_{i+1} + \langle a_{i+1}, a'_i \rangle\, a'_{i+1}
\otimes a_i$ et la matrice $2 \times 2$ des accouplements de $(a_i, a_{i+1})$
avec $(a'_i, a'_{i+1})$.

\subsection*{Le lemme de relèvement (page 94)}

\emph{Lemme.} Soit $\mathcal{E}$ un fibré vectoriel de rang $N \geqslant 3$, $X =
\mathbb{P}(\mathcal{E})$. a) Toute pseudo-réflexion $\tau_X$ de $X$ provient
d'une unique pseudo-réflexion $\tau_\mathcal{E}$ de $\mathcal{E}$, propre si
$\tau_X$ l'est. b) Si $C = \mathbb{P}(V)$ est un sous-fibré projectif et
$\tau_X$ propre de centre dans $C$, $\tau_\mathcal{E}$ induit l'identité sur
$\mathcal{E}/V$ ; si $C$ est un hyperplan et $\mathcal{E}$ l'enveloppe
vectorielle canonique de $X \smallsetminus C$, l'action de $\tau_X$ sur $X
\smallsetminus C$ s'identifie à celle de $\tau_\mathcal{E}$ sur $\pi^{-1}(1)$.

Le lemme est énoncé sans preuve ; l'unicité vient de ce que, en rang $\geqslant
3$, l'hyperplan fixé point par point est déterminé par l'automorphisme
projectif, et que le relèvement qui y vaut l'identité est unique. Suit un
programme : « Dégager » les conditions de stabilité d'un sous-fibré projectif par
des pseudo-réflexions, les conditions de commutation de deux pseudo-réflexions,
l'ordre d'une pseudo-réflexion et celui du produit de deux.

\pagerange{96}{105}
\section*{XVI. Groupes engendrés par des pseudo-réflexions vectorielles (pages 95 à 105)}

\subsection*{Le théorème (pages 96 à 100)}

Dans cette section « libre » s'entend au sens fort : une famille est libre si
elle est la base d'un sous-module facteur direct, comme le précise la page
103\footnote{Sur un corps, c'est l'indépendance linéaire. Sur un anneau,
c'est ce sens fort qui rend (9) équivalente à (12) et (13) ; la page ne le dit
qu'à la page 103 (« strictement libre »).}. Les pseudo-réflexions sont
strictes : $a$ et $a'$ partout non nuls.

\emph{Théorème.} Soient $k$ un anneau, $\mathcal{E}$ un $k$-module, $\tau_i =
\mathrm{id} + a'_i \otimes a_i$ ($0 \leqslant i \leqslant n$) des
pseudo-réflexions, $s_0 \in \mathcal{E}$, $s_{i+1} = \tau_i s_i$,
\[
\mathcal{E}_i = k(s_0, a_0, \ldots, a_{i-1}), \qquad \nu_i = \langle s_0, a'_0
\rangle \langle a_0, a'_1 \rangle \cdots \langle a_{i-1}, a'_i \rangle
\]
\footnote{La page note $\lambda_i$ ces produits ; ce ne sont pas les valeurs
propres des sections précédentes, d'où la lettre $\nu_i$.}.
\begin{enumerate}
  \item[1°)] (5) $s_i \in \mathcal{E}_i$ pour tout $i$ ; et (6) si $a'_j$ est nulle
    sur $\mathcal{E}_{j-1}$ pour $j \leqslant i$, alors $s_{i+1} \equiv \nu_i
    a_i \bmod \mathcal{E}_i$.
  \item[2°)] Les conditions (A) : (7) $\tau_i\tau_j = \tau_j\tau_i$ pour $i, j$ non
    consécutifs ; (8) $\tau_i s_0 = s_0$ pour $i \geqslant 1$ ; (9) $(s_i)_{0
    \leqslant i \leqslant n+1}$ libre ; sont équivalentes aux conditions (B) :
    (10) $\langle a_i, a'_j \rangle = 0$ pour $i, j$ non consécutifs ; (11)
    $\langle s_0, a'_i \rangle = 0$ pour $i \geqslant 1$ ; (12) $(s_0, a_0,
    \ldots, a_n)$ libre ; (13) $\nu_n \in k^*$.
  \item[3°)] Elles entraînent (14) $\mathcal{E}_i = k(s_0, \ldots, s_i)$ ; (15)
    $\tau_j \mathcal{E}_i = \mathcal{E}_i$ pour $j \neq i$ ; (16) $\tau_i
    \mathcal{E}_i \neq \mathcal{E}_i$ si $k \neq 0$\footnote{La page écrit (14)
    $\mathcal{E}_i = k(s_0, \ldots, s_{i-1})$ ; $\mathcal{E}_i$ a $i+1$
    générateurs, et la page 98 écrit $\mathcal{E}_{n+1} = k(s_0, \ldots,
    s_{n+1})$ : glissement d'indice.}.
\end{enumerate}
Les $\mathcal{E}_i$ sont donc le drapeau vectoriel qui relève le drapeau $X_*$ de
la section XV.

\emph{Démonstration.} 1°) Par récurrence, avec $a_{-1} = s_0$, $\nu_{-1} = 1$ :
$s_{i+1} = s_i + \langle s_i, a'_i \rangle a_i \in \mathcal{E}_i + ka_i =
\mathcal{E}_{i+1}$ ; et si $s_i \equiv \nu_{i-1} a_{i-1} \bmod \mathcal{E}_{i-1}$
et $a'_i|_{\mathcal{E}_{i-1}} = 0$, alors $\langle s_i, a'_i \rangle =
\nu_{i-1}\langle a_{i-1}, a'_i \rangle = \nu_i$.

2°) Rappelons que $\tau_i\tau_j(x) = x + \langle x, a'_j \rangle a_j + \langle x,
a'_i \rangle a_i + \langle x, a'_j \rangle\langle a_j, a'_i \rangle a_i$ ; la
commutation équivaut donc à $\langle x, a'_j \rangle\langle a_j, a'_i \rangle
a_i = \langle x, a'_i \rangle\langle a_i, a'_j \rangle a_j$ pour tout $x$.
(B) $\Rightarrow$ (A) : (10) donne (7), (11) donne (8) ; (10) et (11) disent que
$a'_j$ est nulle sur $\mathcal{E}_{j-1}$, donc, par 1°), la matrice qui exprime
$(s_0, \ldots, s_{n+1})$ dans $(s_0, a_0, \ldots, a_n)$ est triangulaire de
diagonale $(1, \nu_0, \ldots, \nu_n)$ ; comme $\nu_n$ est inversible, tous les
$\nu_i$ le sont, et (12) donne (9). (A) $\Rightarrow$ (B), par récurrence sur
$n$ (pages 98 et 99) : on suppose tout établi pour $(s_0, \tau_0, \ldots,
\tau_{n-1})$, avec $\mathcal{E}_n = k(s_0, \ldots, s_n)$. Comme $s_{n+1} = s_n +
\langle s_n, a'_n \rangle a_n$ avec $s_n \in \mathcal{E}_n$, la liberté de
$s_{n+1}$ modulo $\mathcal{E}_n$ entraîne celle de $a_n$, donc (12). Alors $(a_i,
a_n)$ est libre pour $i \leqslant n-2$, et la commutation, avec $a'_i$ et $a'_n$
partout non nulles, donne $\langle a_i, a'_n \rangle = \langle a_n, a'_i \rangle
= 0$ ; de même (8) donne $\langle s_0, a'_n \rangle = 0$. Enfin $s_n \equiv
\nu_{n-1}a_{n-1} \bmod \mathcal{E}_{n-1}$ et $a'_n$ nulle sur $\mathcal{E}_{n-1}$
donnent $\langle s_n, a'_n \rangle = \nu_n$, et la liberté de $s_{n+1}$ modulo
$\mathcal{E}_n$ équivaut à $\nu_n \in k^*$ ; d'où (13) et $\mathcal{E}_{n+1} =
k(s_0, \ldots, s_{n+1})$. « cqfd ».

3°) (15) : pour $j < i$, $a_j \in \mathcal{E}_i$ ; pour $j > i$, $a'_j$ est
nulle sur $\mathcal{E}_i$, donc $\tau_j$ y est l'identité. (16) : $\tau_i$ est
l'identité sur $\mathcal{E}_{i-1}$ et $\tau_i a_{i-1} = a_{i-1} + \langle a_{i-1},
a'_i \rangle a_i$ avec $\langle a_{i-1}, a'_i \rangle$ inversible, donc $\tau_i
\mathcal{E}_i = \mathcal{E}_{i-1} + k(a_{i-1} + \langle a_{i-1}, a'_i \rangle
a_i) \neq \mathcal{E}_i$, $a_i$ étant libre modulo $\mathcal{E}_i$\footnote{La
page écrit $\tau_i\mathcal{E}_i = \mathcal{E}_{i-1} + ka_i$, en omettant le terme
$a_{i-1}$ ; la conclusion est la même.}.

\emph{Corollaire} (pages 97 et 98). Posons $h_i = ka_i$, $H_i = \operatorname{Ker}
a'_i$, $h_{-1} = ks_0$. Les conditions (A') : (15) $h_i \subset H_j$ pour $-1
\leqslant i < j \leqslant n$, $j \neq i+1$, et (15') $(s_i)$ libre --- ou encore
les $h_i$ indépendants et $h_i \not\subset H_{i+1}$ sur chaque fibre ---, et
(B') : (17) $\langle a_i, a'_j \rangle = 0$ pour $-1 \leqslant i < j \leqslant
n$, $j \neq i+1$, (18) $(a_i)_{-1 \leqslant i \leqslant n}$ libre, (19) $\nu_n
\in k^*$, sont équivalentes et entraînent (14). (17) n'est que la traduction de
(15) ; le reste est l'argument de 2°)\footnote{Ces conditions ne portent que sur
les $i < j$ : c'est la moitié « triangulaire » des conditions de commutation,
celle qui suffit à construire le drapeau. Le numéro (15) sert deux fois sur la
page.}.

\emph{Corollaire 2} (page 100). Sous les conditions du théorème, avec $H_i =
\mathcal{E}^{\tau_i} = \operatorname{Ker} a'_i$ et $h_i = \operatorname{Im}(\tau_i
- \mathrm{id}) = ka_i$ (21), on a : (22) $\mathcal{E}_i \subset H_j$ pour $i <
j$ ; (23) $\mathcal{E}_i \cap H_i = \mathcal{E}_{i-1}$ ; (24) $\tau_i\tau_{i+1}
\neq \tau_{i+1}\tau_i$, même modulo $k^*$, si $k \neq 0$ ; (25) $\langle a_i,
a'_{i+1} \rangle \in k^*$ pour $-1 \leqslant i \leqslant n-1$ ; (25') $h_i
\not\subset H_{i+1}$ si $k \neq 0$. (21) vient de la liberté de $a'_i$ et de sa
surjectivité, $\langle a_{i-1}, a'_i \rangle$ étant inversible ; (25) vient de
(13), $\nu_n$ étant le produit de ces accouplements ; (23) de ce que $a'_i$ est
nulle sur $a_{-1}, \ldots, a_{i-2}$ et inversible sur $a_{i-1}$.

Pour (24) (page 101) : si $\tau_i\tau_{i+1} = r\,\tau_{i+1}\tau_i$ avec $r \in
k^*$, l'égalité sur $H_i \cap H_{i+1}$, non nul en rang $\geqslant 3$, force $r
= 1$ ; et le commutateur vaut
\[
\tau_i\tau_{i+1} - \tau_{i+1}\tau_i = \beta_i\, a'_{i+1} \otimes a_i - \langle
a_i, a'_{i+1} \rangle\, a'_i \otimes a_{i+1},
\]
qui envoie $a_{i-1}$ sur $-\langle a_{i-1}, a'_i \rangle\langle a_i, a'_{i+1}
\rangle\, a_{i+1} \neq 0$, puisque $\langle a_{i-1}, a'_{i+1} \rangle = 0$. Donc
$\tau_i$ et $\tau_{i+1}$ ne commutent sur aucune fibre, quel que soit
$\beta_i$\footnote{La page écrit le commutateur à $r$ près sans le terme en
$(1-r)$, et marque « $\in k^*$ » sous $\langle a_{i+1}, a'_i \rangle$, qui est
$\beta_i$ et n'est pas supposé inversible ; elle conclut par la surjectivité de
$(a'_i, a'_{i+1}) \colon \mathcal{E} \to k^2$, de matrice triangulaire à
diagonale inversible. Ni l'inversibilité de $\beta_i$ ni toute la surjectivité ne
sont nécessaires : l'évaluation en $a_{i-1}$ suffit. On a rédigé l'argument
ainsi.}. Enfin (25') suit de (25).

\subsection*{L'existence du point de départ (pages 102 à 105)}

\emph{Remarque} (page 102). Si $\mathcal{E}$ est localement libre de rang $n+2$,
on obtient des conditions équivalentes $(\bar A)$, $(\bar B)$ en renforçant (9) en
« $(s_i)_{0 \leqslant i \leqslant n+1}$ est une base de $\mathcal{E}$ » et (12) en
« $(s_0, a_0, \ldots, a_n)$ est une base » ; elles équivalent aussi à (A), (B) et
(26) $\mathcal{E}_{n+1} = \mathcal{E}$. De même pour le corollaire\footnote{La
page écrit (9') « $(a_i)_{0 \leqslant i \leqslant n+1}$ forment une base de $E$ »
et (12') « $(a'_i)_{-1 \leqslant i \leqslant n}$ forment une base de $E$ » ; les
indices et le module (les $a'_i$ sont des formes) indiquent la lecture
ci-dessus, qui est celle de (26).}.

\emph{Proposition} (pages 102 et 103). Soient $\mathcal{E}$ localement libre de rang
$n+2$ et $z_i = \mathrm{id} + a'_i \otimes a_i$ ($0 \leqslant i \leqslant n$) des
pseudo-réflexions. Pour qu'il existe localement $s_0 \in \mathcal{E}$
satisfaisant $(\bar A)$, $(\bar B)$ --- conditions de commutation et de
non-dégénérescence ---, il faut et il suffit que : a) $(a_i)_{0 \leqslant i
\leqslant n}$ soit strictement libre ; a') $(a'_i)$ soit strictement libre dans
$\mathcal{E}^\vee$ ; b) $\langle a_i, a'_j \rangle = 0$ pour $i < j$, $j \neq
i+1$ ; c) $\langle a_i, a'_{i+1} \rangle \in k^*$ ($0 \leqslant i \leqslant
n-1$), c'est-à-dire $h_i \not\subset H_{i+1}$ sur chaque fibre.

\emph{Démonstration.} La nécessité de a), b), c) est claire. Pour a'), soit
$a'_{n+1}$ la forme de noyau $k(a_{-1}, \ldots, a_{n-1})$ ; la matrice des
$\langle a_i, a'_{j} \rangle$ ($-1 \leqslant i \leqslant n$, $0 \leqslant j
\leqslant n+1$) est triangulaire, de diagonale $c_{-1} = \langle s_0, a'_0
\rangle, c_0 = \langle a_0, a'_1 \rangle, \ldots, c_n = \langle a_n, a'_{n+1}
\rangle$, et son déterminant $\nu_n c_n$ est inversible ; donc $(a'_0, \ldots,
a'_{n+1})$ est une base de $\mathcal{E}^\vee$. Inversement (pages 104 et 105),
sous a), a'), b), c), les $s$ tels que $\langle s, a'_i \rangle = 0$ pour $i
\geqslant 1$ forment, par a'), un facteur direct $F$ de rang $2$. Dans $F$,
la condition (19) devient $\langle s_0, a'_0 \rangle \in k^*$, c'est-à-dire
$s_0 \notin D = F \cap \operatorname{Ker} a'_0$, une droite. Reste (18), la
liberté de $(s_0, a_0, \ldots, a_n)$ : soit $L = \mathcal{E}/k(a_0, \ldots,
a_n)$, de rang $1$ ; $F \to L$ est surjectif, car sur un corps $F \cap k(a_0,
\ldots, a_{n-1}) = 0$ --- les formes $a'_1, \ldots, a'_n$ ont sur $k(a_0, \ldots,
a_{n-1})$ une matrice triangulaire de diagonale $c_0, \ldots, c_{n-1}$
inversible --- et $k(a_0, \ldots, a_{n-1})$ est de codimension $1$ dans $k(a_0,
\ldots, a_n)$ alors que $F$ est de dimension $2$. Avec $\Delta =
\operatorname{Ker}(F \to L)$, (18) dit $s_0 \notin \Delta$. Donc les $s_0$
cherchés sont les sections du fibré $F$ de rang $2$ qui ne sont, sur aucune
fibre, ni dans $D$ ni dans $\Delta$ : il en existe toujours localement --- même
sur un corps à deux éléments, un plan ayant trois droites\footnote{Cette
précision est de cette lecture ; la page dit « il en existe toujours localement,
ok ».}.

\pagerange{107}{112}
\section*{XVII. Classes de Stiefel-Whitney et classe de Brauer de la forme invariante (pages 106 à 112)}

On suppose $2$, les $\beta_i$ et $\delta(f_V) = B\rho$ inversibles ; $\partial
\colon H^0(-, \mathbb{G}_m) \to H^1(-, \mu_2) = H^1(-, \mathbb{Z}/2)$ est le
cobord de Kummer, et $w_*$ la classe de Stiefel-Whitney totale d'une forme
quadratique, de sorte que $w_*(\langle d_1, \ldots, d_r \rangle) = \prod (1 +
\partial d_i)$\footnote{Ce sont les classes de Stiefel-Whitney des formes
quadratiques au sens de Delzant (1962), que Grothendieck pouvait connaître ;
la page n'en cite pas la source.}.

\emph{La décomposition} (page 107). L'orthogonal dans $V$ de $\mathcal{O}_S(a_1,
\ldots, a_n)$ est, puisque $\psi(a_i) = -c_i a'_i$, le noyau commun des $a'_i$
($i \geqslant 1$) dans $V$ : il est engendré par
\[
w = \tilde c - \rho s_0 = -\rho\,\tilde s_0 = \sum_{i \geqslant 0} \xi_i a_i, \qquad
f_V(w) = \rho^2 f_V(\tilde s_0) = \tilde\rho\tilde\sigma ,
\]
où $\tilde s_0 = s_0 - c$ est la projection orthogonale de $s_0$ sur $V$ ; il
retrouve $f_V(\tilde s_0) = \tilde\rho\tilde\sigma/\rho^2$ comme à la page 74.
Or la forme induite sur $\mathcal{O}_S(a_1, \ldots, a_n)$ est $\beta_0$ fois la
forme $f_V(\beta_1, \ldots, \beta_{n-1})$ du polytope de rang inférieur, puisque
$f(a_i) = \beta_0 \cdot (\beta_1 \cdots \beta_{i-1})$. Si de plus $\sigma$ est
inversible --- le réseau $\mathcal{O}_S w \oplus \mathcal{O}_S(a_1, \ldots,
a_n)$ est d'indice $\xi_0 = \sigma$ dans $V$ ---, on a donc la décomposition
orthogonale
\[
f_V(\beta_0, \ldots, \beta_{n-1}) \simeq \langle \tilde\rho\tilde\sigma \rangle
\perp \beta_0 \cdot f_V(\beta_1, \ldots, \beta_{n-1}) .
\]
En l'itérant, avec $\tilde\rho_i = \tilde\rho_{i,n-1} = L'_{n-i+1}(\beta_i,
\ldots, \beta_{n-1})$ ($\tilde\rho_0 = \tilde\rho$, $\tilde\rho_1 = \tilde\sigma$,
$\tilde\rho_n = \tilde\rho_{n+1} = 1$), on trouve, si tous les $\tilde\rho_i$ sont
inversibles,
\[
f_V \simeq \mathop{\perp}_{i=0}^{n} \bigl\langle (\beta_0 \cdots
\beta_{i-1})\, \tilde\rho_i\,\tilde\rho_{i+1} \bigr\rangle ,
\qquad
w_*(f_V) = \prod_{i=0}^{n} \bigl(1 + \partial(\beta_0 \cdots \beta_{i-1}) +
\eta_i + \eta_{i+1}\bigr), \quad \eta_i = \partial\tilde\rho_i .
\]
Le produit des coefficients est $B\tilde\rho$ à un carré près, comme il se
doit\footnote{La décomposition et la formule sont de cette lecture. La page écrit
la récurrence $w_*(f_{V_{n+1}}(\beta_0, \ldots)) = [1 + \partial(\tilde\rho_0
\tilde\rho_1)]\, w_*(f_{V_n}(\beta_1, \ldots))$, qui omet le facteur
$\beta_0$ de la forme induite, et en tire (page 109) $w_*(f_V) = \prod_{i=0}^{n-1}
(1 + \eta_i + \eta_{i+1})$ avec $\eta_n = \partial\beta_{n-1}$. Cette formule est
fausse en général : pour $n = 2$, $\beta_0 = -1$, $\beta_1 = 1$ sur
$\mathbb{R}$, on a $\tilde\rho = 4$, $\tilde\sigma = 3$, $f_V = x_0^2 - x_1^2 -
x_2^2 + x_0x_1 + x_1x_2$ est de signature $(1, 2)$ et $w_2(f_V) =
\partial(-1)^2 \neq 0$, alors que la formule de la page donne $w_* = 1$. Pour
$n = 1$ elle est juste. Le calcul de rang $2$ de la page 108, $f(\beta a_0 +
2a_1) = \beta(4 - \beta)$, d'où $w_*(f_V(\beta)) = 1 + \partial(\beta(4 -
\beta))$, est juste.}. La décomposition est valable sous l'hypothèse que tous
les $\tilde\rho_i$ sont inversibles, hypothèse qu'il juge lui-même « parasite »
page 110. Comme $f_V(a_0) = 1$, on a $f_V \simeq \langle 1 \rangle \perp q'$,
donc
\[
w_{n+1}(f_V) = 0 ,
\]
ce que dit son « NB » de la page 109. Pour $f_E$, comme $V^\perp =
\mathcal{O}_S\tilde c$ avec $f_E(\tilde c) = -\tilde\rho\tilde\sigma$ (et $\rho$
inversible), $w_*(f_E) = (1 + \varepsilon + \eta_0 + \eta_1)\, w_*(f_V)$, avec
$\varepsilon = \partial(-1)$. Il se demande enfin si, en menant le calcul dans
l'autre sens, on trouverait « d'autres expressions remarquables ».

\emph{Le cas $n = 2$} (pages 110 et 111). Sans l'hypothèse parasite, il prend
pour base orthogonale $a_0$, $a_2$ et $u = \beta_0 a_0 + 2a_1 + a_2$, orthogonal
aux deux, avec $f(a_0) = 1$, $f(a_2) = \beta_0\beta_1$, $f(u) = -\beta_0(\alpha_0 +
\alpha_1) = \beta_0\tilde\rho$ ; d'où, avec
\[
\varepsilon = \partial(-1), \quad \eta_0 = \partial\beta_0, \quad \eta_1 =
\partial\beta_1, \quad \eta_2 = \partial(-(\alpha_0 + \alpha_1)) = \partial\tilde\rho ,
\]
\[
w_*(f_V) = (1 + \eta_0 + \eta_1)(1 + \eta_0 + \eta_2), \qquad w_1 = \eta_1 +
\eta_2, \qquad w_2 = \eta_0^2 + \eta_0\eta_1 + \eta_1\eta_2 + \eta_2\eta_0 .
\]
C'est juste, et s'accorde avec la formule générale ci-dessus grâce à la
relation de Steinberg $\partial a \cdot \partial b = 0$ pour $a + b = 1$,
appliquée à $\tilde\rho/\tilde\sigma + \beta_0/\tilde\sigma = 1$\footnote{Vérification de cette lecture.}. Il pose ensuite
\[
c(f_V) = w_2(f_V) + w_1(f_V)^2 = (\eta_0 + \eta_1 + \eta_2)^2 + (\eta_0\eta_1 +
\eta_1\eta_2 + \eta_2\eta_0),
\]
nul si et seulement si la conique $Q_C$ est isomorphe à celle de $x^2 + y^2 +
z^2$, et
\[
\operatorname{Br}(f_V) = c(f_V) - c(X^2 + YZ) = (\varepsilon + \eta_0 + \eta_1 +
\eta_2)^2 + (\eta_0\eta_1 + \eta_1\eta_2 + \eta_2\eta_0),
\]
puisque $c(X^2 + YZ) = \varepsilon^2$. Pour une forme ternaire $\langle a, b, c
\rangle$, la classe de la conique est celle de l'algèbre de quaternions $(-ab,
-ac)$, et un calcul bilinéaire donne $(-ab, -ac) = w_2 + w_1^2 + (-1, -1)$ :
$\operatorname{Br}(f_V)$ est donc bien, comme il le dit, la classe de Brauer du
fibré en coniques $Q_C$, dont l'annulation signifie que ce fibré est le
projectivisé d'un fibré vectoriel de rang $2$\footnote{C'est aujourd'hui
l'invariant de Clifford, ou de Hasse-Witt selon la normalisation ; la conique
est le schéma de Severi-Brauer de l'algèbre de Clifford paire, comme dans le
dossier 82.}.

\emph{La base universelle} (pages 111 et 112). Le calcul universel se fait sur
\[
S_0 = \operatorname{Spec} \mathbb{Z}[\beta_0, \beta_1][1/2\beta_0\beta_1(4 -
\beta_0 - \beta_1)].
\]
Il affirme que $\varepsilon, \eta_0, \eta_1, \eta_2$ forment
une base de $H^1(S_0, \mathbb{Z}/2)$, et qu'il est « plausible » que les dix
produits $\varepsilon^2, \eta_i^2, \varepsilon\eta_i, \eta_i\eta_j$ soient
linéairement indépendants dans $H^2(S_0, \mathbb{Z}/2)$, avec un point
d'interrogation. Les deux points demandent correction. Les unités de l'anneau
de $S_0$ sont, au signe près, les monômes en $2, \beta_0, \beta_1, \tilde\rho$, et
son groupe de Picard est nul : $H^1(S_0, \mathbb{Z}/2)$ est de dimension $5$, de
base $\varepsilon, \partial 2, \eta_0, \eta_1, \eta_2$. Et comme $S_0$ est
régulier de Picard nul, $H^2(S_0, \mu_2)$ s'injecte dans le groupe de Brauer de
son corps des fractions, où vaut la relation $\partial a \cdot \partial a =
\partial a \cdot \partial(-1)$ ; donc $\eta_i^2 = \varepsilon\eta_i$, et les dix
produits ne sont pas indépendants. L'indépendance des sept qui restent n'est pas
établie ici\footnote{Les deux corrections sont de cette lecture. La page 112 sait
que $\partial 2$ compte : elle le met dans la liste du « NB » qui suit.}.

\emph{Le problème} (page 112) : calculer les classes de Stiefel-Whitney totales de
la forme quadratique lisse « à coefficients indéterminés » sur $\mathbb{Z}[(c_{ij})_{1
\leqslant i \leqslant j \leqslant n}][1/2\delta']$. Le $H^1$ de cette base est
engendré par $\varepsilon = \partial(-1)$, $\partial 2$ et $\eta = \partial\delta' =
w_1(f)$ ; les deux premiers s'annulent sur $\mathbb{Q}(\sqrt{-1}, \sqrt 2)$, a
fortiori sur $\mathbb{C}$, et il conclut qu'on ne peut exprimer
multiplicativement à partir de ces éléments aucune des classes $w_2, w_3,
\ldots$ Le raisonnement suppose que, sur $\mathbb{C}$, $w_2$ n'est pas un
multiple de $w_1^2$, ce que la page ne montre pas ; la conclusion reste un
constat à étayer\footnote{La lettre du corps de base est surchargée (un
$\mathbb{Z}$ barré ?) et lue $\mathbb{Q}$.}.

\end{document}
